\newcommand{\vX}{\mathbf{X}}
\newcommand{\vY}{\mathbf{Y}}

\renewcommand{\paragraph}[1]{\par\medskip\noindent\textbf{#1}\enspace}

\newpage

\subsection{Margin Bounds Setup}
\label{app:theory:margin_bounds}
We lay out the setup we'll use throughout the appendix to prove our margin bounds.
\subsection*{Stored items and kernel}
We store $F$ key--value items $\{(\vk_{t},\vv_{t})\}_{t=1}^F$ with keys $\vk_{t}\in\R^d$ and values $\vv_{t}\in\R^d$.
A feature map $\phi:\R^d\to\R^p$ induces a kernel $K:\R^d\times \R^d \to \R$
\[
\hat{\matK}_{ti} \defeq K(\vk_{t}, \vk_{i})=\ip{\phi(\vk_{t})}{\phi(\vk_{i})}\qquad (t,i\in[F]).
\]
We write $\hat{\matK}\in\R^{F\times F}$ for the Gram matrix.

\subsection*{Codes and retrieval output}
Each index $t\in[F]$ has an associated code vector $\vc_{t}\in\R^{d}$.
Given a stored query at index $i$ (i.e. query key $\vk_{i}$), the retrieval output is
\begin{equation}
 \vy_{i} \defeq \sum_{t=1}^F \vc_{t}\,\hat{\matK}_{ti}\in\R^{d}.
 \label{eq:yi}
\end{equation}
Unless otherwise noted, for simplicity, we assume $\vc_{t}=\vv_{t}$ for all the upcoming theorems and their proofs.

\subsection*{Pairwise margin}
For a stored index $i$ and a competitor $j\neq i$, we define the pairwise margin
\begin{equation}
\gamma_{ij} \defeq \ip{\vv_{i}-\vv_j}{\vy_i}.
\label{eq:margin}
\end{equation}
We write $\gamma_{\min}\defeq\min_{i\neq j}\gamma_{ij}$ for the worst-case (minimum) pairwise margin.
Expanding, we see that the pairwise margin decomposes into signal and cross-talk components.

\begin{equation}
\gamma_{ij}
= \underbrace{\hat{\matK}_{ii}\,\ip{\vv_{i}-\vv_j}{\vc_i}}_{\text{signal }}
+ \underbrace{\sum_{t\neq i}\hat{\matK}_{ti}\,\ip{\vv_{i}-\vv_j}{\vc_t}}_{\text{cross-talk }}.
\label{eq:signal-crosstalk}
\end{equation}

\subsection*{Margin convention (absorbing the competitor)}
For the pairwise margin $\gamma_{ij}$, the single term $t=j$ inside the cross-talk sum in \eqref{eq:signal-crosstalk}
corresponds to the specific \emph{competitor} item. For simplicity of our proofs, in the rest of the appendix, we re-define the signal and cross-talk by absorbing the competitor term into the signal:

\begin{equation}
\gamma_{ij} = \underbrace{\widetilde{s}_{ij}}_{\text{signal}} + \underbrace{\widetilde{z}_{ij}}_{\text{cross-talk}}
\label{eq:new-signal-crosstalk}
\end{equation}
with
\begin{equation}
    \widetilde{s}_{ij} \defeq \hat{\matK}_{ii}\,\ip{\vv_{i}-\vv_j}{\vc_i} + \hat{\matK}_{ji}\,\ip{\vv_{i}-\vv_j}{\vc_j}, \qquad
    \widetilde{z}_{ij} \defeq \sum_{t\notin\{i,j\}}\hat{\matK}_{ti}\,\ip{\vv_{i}-\vv_j}{\vc_t}
    \label{eq:new-signal-crosstalk-sij-zij}
\end{equation}

\subsection{Embedding Geometric Summary Statistics Definitions}\label{app:emb-geo-defns}
We collect here the formal definitions of all geometric summary statistics that enter our margin bounds.
Throughout, $\hat{\matK}$ denotes the kernel Gram matrix with entries $\hat{\matK}_{ti} = K(\vk_t, \vk_i)$, and $\vv_1,\dots,\vv_F \in \R^d$ are the stored value vectors with associated Hebbian codes $\vc_1,\dots,\vc_F \in \R^d$.

\paragraph{Key-geometry statistics.}

\begin{definition}[Diagonal and off-diagonal kernel energies]\label{def:Kdiag}
\[
K_{\min}^{\mathrm{diag}} \;\defeq\; \min_{i \in [F]}\, \hat{\matK}_{ii},
\qquad
K_{\max}^{\mathrm{off}} \;\defeq\; \max_{i \neq j}\, |\hat{\matK}_{ij}|.
\]
$K_{\min}^{\mathrm{diag}}$ is the minimum kernel self-similarity; $K_{\max}^{\mathrm{off}}$ is the maximum off-diagonal kernel entry.
\end{definition}

\begin{definition}[Kernel column energy]\label{def:Ecol}
\[
E_K \;\defeq\; \max_{i \in [F]}\, \sum_{t \neq i} \hat{\matK}_{ti}^2.
\]
$E_K$ measures the worst-case squared $\ell_2$ energy of an off-diagonal kernel column, capturing key-embedding crowding: it grows when keys cluster in embedding space and kernel overlaps are large.
\end{definition}

\paragraph{Value-geometry statistics.}

\begin{definition}[Value separability]\label{def:Vmin}
\[
V_{\min} \;\defeq\; \min_{i \neq j}\, \langle \vv_i - \vv_j,\, \vv_i \rangle,
\qquad
V_{\max} \;\defeq\; \max_{i \neq j}\, |\langle \vv_i - \vv_j,\, \vv_j \rangle|.
\]
$V_{\min}$ is the signal-side value separability floor; $V_{\max}$ is the cross-talk-side value interaction ceiling.
\end{definition}

For each pair $i \neq j$, let $\vY^{(ij)} \defeq \bigl(\langle \vv_i - \vv_j,\, \vc_t \rangle\bigr)_{t \notin \{i,j\}} \in \R^{F-2}$ denote the vector of value-difference inner products with all non-target, non-competitor codes, and let $\vone \in \R^{F-2}$ denote the all-ones vector.

\begin{definition}[Mean competitor alignment]\label{def:BY}
\[
B_Y
\defeq
\max_{i\neq j}
\left|
\sum_{t\notin\{i,j\}}
\left\langle \vv_i-\vv_j,\vc_t\right\rangle
\right|
=
\max_{i\neq j}
\left|
\left\langle \vone,\vY^{(ij)}\right\rangle
\right|.
\]
$B_Y$ controls the bias contribution to cross-talk arising from the mean of the kernel off-diagonal entries (equal to $1/d$ under isotropic keys).
\end{definition}



\begin{definition}[Value-difference energy]\label{def:EY}
\[
E_v \;\defeq\; \max_{i \in [F]}\max_{j \neq i}\, \sum_{t \neq i} \langle \vv_i - \vv_j,\, \vc_t \rangle^2.
\]
$E_v$ measures the worst-case squared $\ell_2$ energy of the value-difference inner-product vector, capturing value-embedding interference: it grows when value embeddings cluster and non-target codes align with the target direction.
\end{definition}

\begin{definition}[Value sparsity]\label{def:Lv}
Recall $\vY^{(ij)}$ from \cref{def:BY}. Define
\[
L_v \;\defeq\; \max_{i \neq j}\, \frac{\|\vY^{(ij)}\|_1^2}{\|\vY^{(ij)}\|_2^2}
\]
$L_v \in [1, F-2]$ is an effective sparsity parameter: it equals $1$ when the energy of $\vY^{(ij)}$ is concentrated on a single coordinate and equals $F-2$ when it is spread uniformly. It amplifies cross-talk when value-difference inner products have heavy-tailed distributions across competitors.
\end{definition}

\paragraph{Key--value coupling.}

\begin{definition}[Coupling factor]\label{def:kappa}
For stored index $i$ and competitor $j \neq i$, let $E_K(i) \defeq \sum_{t \neq i} \hat{\matK}_{ti}^2$ and $E_v(i,j) \defeq \sum_{t \neq i} \langle \vv_i - \vv_j, \vc_t \rangle^2$ be the pairwise energies. Define the pairwise coupling
\[
\kappa^{ij}
\;\defeq\;
\frac{\displaystyle\left|\sum_{t \neq i} \hat{\matK}_{ti}\,\langle \vv_i - \vv_j,\, \vc_t \rangle\right|}{\sqrt{E_K(i)}\,\sqrt{E_v(i,j)}}
\;\in\; [0,1],
\]
\[
\kappa \;\defeq\; \max_{i \in [F]}\,\max_{j \neq i}\; \kappa^{ij}.
\]
$\kappa$ quantifies the worst-case alignment between the kernel column pattern and the value-interference pattern: $\kappa = 0$ when the two are orthogonal and $\kappa = 1$ when they are perfectly aligned.
\end{definition}

\subsection{Cross-Talk Bounds}
We start by providing upper bounds for the cross-talk term $\widetilde{z}_{ij}$.
\paragraph{Cross-talk as inner product.} First, we define cross-talk as an inner product. Concretely, fix $(i,j)$ with $j\neq i$ and define the off-diagonal kernel column
\begin{equation}
\vX^{(ij)} \defeq \big(\hat{\matK}_{ti}\big)_{t\notin \{i,j\}}\in\R^{F-2}.
\label{eq:Xdef}
\end{equation}
And for the value/code side:
\begin{equation}
\vY^{(ij)} \defeq \big(\ip{\vv_{i}-\vv_j}{\vc_t})_{t\notin \{ i,j\}}\in\R^{F-2},
\label{eq:Ydef}
\end{equation}
Then the cross-talk term is the inner product
\begin{equation}
 \widetilde{z}_{ij}
 \defeq \sum_{t\notin\{i,j\}}\hat{\matK}_{ti}\,\ip{\vv_{i}-\vv_j}{\vc_t}
 =\ip{\vX^{(ij)}}{\vY^{(ij)}}.
 \label{eq:Zij}
\end{equation}
We consider the largest (worst case) possible cross-talk:
\begin{equation}
 \widetilde{z}_{\max}\ \defeq\ \max_{i\neq j}|\widetilde{z}_{ij}|.
 \label{eq:zmax}
\end{equation}
Note that $\widetilde{z}_{ij}\le \widetilde{z}_{\max}$ for all $i\neq j$.


\paragraph{Cross-talk summary statistics.} For each $(i,j)$ define the (squared) energies
\begin{equation}
\widetilde{E}_{v}(i,j) \defeq \norm{\vY^{(ij)}}_{2}^2 = \sum_{t\notin\{i,j\}}\ip{\vv_{i}-\vv_j}{\vc_t}^2,
\qquad
\widetilde{E}_{K}(i,j) \defeq \norm{\vX^{(ij)}}_{2}^2 = \sum_{t\notin\{i,j\}}\hat{\matK}_{ti}^2.
\label{eq:EvEk}
\end{equation}
Define the coupling factor
\begin{equation}
\widetilde{\kappa}^{ij} \defeq \frac{|\ip{\vX^{(ij)}}{\vY^{(ij)}}|}{\norm{\vX^{(ij)}}_{2}\,\norm{\vY^{(ij)}}_2}\in[0,1],
\label{eq:kappaij}
\end{equation}
(setting $\widetilde{\kappa}^{ij}=0$ if a norm is zero).
Now define the worst-case summary statistics
\begin{equation}
\widetilde{E}_{v} \defeq \max_{i\neq j} \widetilde{E}_{v}(i,j),
\qquad
\widetilde{E}_{K} \defeq \max_{i\neq j} \widetilde{E}_{K}(i,j),
\qquad
\widetilde{\kappa} \defeq \max_{i\neq j} \widetilde{\kappa}^{ij}.
\label{eq:summary-stats}
\end{equation}

\subsubsection{Arbitrary Keys, Arbitrary Values}
\label{app:deterministic}

\begin{theorem}[Cross-talk bound --- arbitrary keys, arbitrary values]
\label{lem:deterministic}
With the summary statistics (\eqref{eq:summary-stats}),
\begin{equation}
\widetilde{z}_{\max}\ \le\ \sqrt{\widetilde{E}_K}\,\sqrt{\widetilde{E}_v}\,\widetilde{\kappa}.
\label{eq:deterministic-bound-app}
\end{equation}
\end{theorem}

\begin{proof}
Fix $i\in[F]$ and $j\neq i$.
If $\|\vX^{(ij)}\|_{2}=0$ or $\|\vY^{(ij)}\|_{2}=0$, then $\widetilde{z}_{ij}=0$ and $\widetilde{\kappa}^{ij}=0$, so the claim is trivial.
Otherwise, by definition of $\widetilde{\kappa}^{ij}$ in \eqref{eq:kappaij},
\[
|\widetilde{z}_{ij}| = |\ip{\vX^{(ij)}}{\vY^{(ij)}}| = \|\vX^{(ij)}\|_{2}\,\|\vY^{(ij)}\|_{2}\,\widetilde{\kappa}^{ij}
= \sqrt{\widetilde{E}_{K}(i,j)}\,\sqrt{\widetilde{E}_{v}(i,j)}\,\widetilde{\kappa}^{ij}.
\]
Using $\widetilde{E}_{K}(i,j)\le \widetilde{E}_{K}$, $\widetilde{E}_{v}(i,j)\le \widetilde{E}_{v}$, and $\widetilde{\kappa}^{ij}\le\widetilde{\kappa}$ and taking $\max_{i\neq j}$ yields \eqref{eq:deterministic-bound-app}.
\end{proof}


\subsubsection{Arbitrary Keys, Isotropic Values}
\begin{theorem}[Cross-talk bound --- arbitrary keys, isotropic values]
\label{app:regime1}
Assume $\vv_{1},\dots,\vv_{F}$ are i.i.d.\ uniform on $\Sph\subset\R^{d}$.
Treat the kernel matrix $\hat{\matK}$ as arbitrary.
Fix $\delta\in(0,1)$ and let
\[
L\ \defeq\ \log\Big(\frac{C_{0} F^2}{\delta}\Big)
\]
for a sufficiently large absolute constant $C_{0}$.
Then with probability at least $1-\delta$ (over the values), the following hold:
\begin{align*}
\widetilde{E}_{v}
&\ \le\ C_{1}\,\frac{(F-2)+L}{d-1},\\
\widetilde{\kappa}
&\ \le\ C_{2}\,\sqrt{\frac{L}{F-2}}.
\end{align*}
Consequently, combining these summary-statistic bounds with Theorem~\ref{lem:deterministic} yields
\[
\widetilde{z}_{\max}
\ \le\
 C_{3}\,\sqrt{\widetilde{E}_K}\,\sqrt{\frac{L}{d-1}}\,\sqrt{1+\frac{L}{F-2}}.
\]
In particular, if $F-2\ge L$ and $d\ge 2$, then
\[
\widetilde{z}_{\max}
\ \le\
 C_{4}\,\sqrt{\widetilde{E}_K}\,\sqrt{\frac{L}{d}}.
\]
\end{theorem}

\begin{proof}
The proof combines two auxiliary concentration bounds via a union bound.

\paragraph{Concentration of $\widetilde{E}_{v}$ (invoke Lemma~\ref{lem:Ev-iso-values}).}
Apply Lemma~\ref{lem:Ev-iso-values} with failure probability $\delta/2$.
This yields an event $\mathcal{E}_{1}$ such that
$\Pbb(\mathcal{E}_{1}^c)\le \delta/2$ and on $\mathcal{E}_{1}$,
\(
\widetilde{E}_{v} \le C\frac{(F-2)+\widetilde L}{d-1}
\)
with $\widetilde L=\log(\tfrac{C_{0} F^2}{\delta/2})$.
Since $\widetilde L=\log(\tfrac{2C_{0} F^2}{\delta})\le \log(\tfrac{C_{0}' F^2}{\delta})$ for $C_{0}'=2C_{0}$,
we may rewrite the bound using $L=\log(\tfrac{C_{0}' F^2}{\delta})$ after adjusting the leading constant.

\paragraph{Concentration of $\widetilde{\kappa}$ (invoke Lemma~\ref{lem:kappa-iso-values}).}
Apply Lemma~\ref{lem:kappa-iso-values} with failure probability $\delta/2$.
This yields an event $\mathcal{E}_{2}$ such that
$\Pbb(\mathcal{E}_{2}^c)\le \delta/2$ and on $\mathcal{E}_{2}$,
\(
\widetilde{\kappa} \le C\sqrt{\tfrac{\widetilde L}{F-2}},
\)
with the same logarithmic factor $\widetilde L$ as above.
Again we rewrite in terms of $L=\log(\tfrac{C_{0} F^2}{\delta})$ by enlarging the absolute constant.

\paragraph{Union bound and plug-in.}
By a union bound,
\[
\Pbb(\mathcal{E}_{1}\cap\mathcal{E}_{2})\ \ge\ 1-\Pbb(\mathcal{E}_{1}^c)-\Pbb(\mathcal{E}_{2}^c)\ \ge\ 1-\delta.
\]
On $\mathcal{E}_{1}\cap\mathcal{E}_{2}$, plug the bounds on $\widetilde{E}_{v}$ and $\widetilde{\kappa}$ into
Theorem~\ref{lem:deterministic} to obtain the stated cross-talk bound.
If additionally $F-2\ge L$, then $\sqrt{1+L/(F-2)}\le \sqrt{2}$ and hence
\[
\widetilde{z}_{\max}\ \le\ C\,\sqrt{\widetilde{E}_K}\,\sqrt{\frac{L}{d-1}}\ \le\ C\,\sqrt{\widetilde{E}_K}\,\sqrt{\frac{L}{d}}
\]
which is exactly the final bound stated in the theorem.
\end{proof}


\subsubsection{Isotropic Keys, Arbitrary Values (Bilinear Kernel)}
\begin{lemma}[Mean$+$residual deterministic cross-talk decomposition]
\label{lem:mean-residual-app}
Fix a pair $(i,j)$ with $j\neq i$.
For any scalar $\mu\in\R$, define the \emph{centered} truncated column
\[
\vX^{\circ(ij)} \defeq \vX^{(ij)}-\mu\,\vone.
\]
Then
\[
\widetilde{z}_{ij}=\ip{\vX^{(ij)}}{\vY^{(ij)}}
= \mu\,\ip{\vone}{\vY^{(ij)}} + \ip{\vX^{\circ(ij)}}{\vY^{(ij)}}.
\]
Consequently, define
\[
B_{Y} \defeq \max_{i\neq j}\big|\ip{\vone}{\vY^{(ij)}}\big|,
\qquad
E_{K}^\circ(i,j) \defeq \|\vX^{\circ(ij)}\|_{2}^2,
\qquad
\kappa^{\circ(ij)} \defeq \frac{|\ip{\vX^{\circ(ij)}}{\vY^{(ij)}}|}{\|\vX^{\circ(ij)}\|_{2}\,\|\vY^{(ij)}\|_2},
\]
with the convention $\kappa^{\circ(ij)}=0$ if $\|\vX^{\circ(ij)}\|_{2}\,\|\vY^{(ij)}\|_{2}=0$.
Also define the worst-case centered energy and centered coupling
\[
E_{K}^\circ \defeq \max_{i\neq j}E_{K}^\circ(i,j),
\qquad
\kappa^\circ \defeq \max_{i\neq j}\kappa^{\circ(ij)}.
\]
Recalling $\widetilde{E}_{v}$ from \eqref{eq:summary-stats}, we have the deterministic bound
\[
\widetilde{z}_{\max}\ \le\ |\mu|\,B_{Y}\ +\ \sqrt{E_{K}^\circ}\,\sqrt{\widetilde{E}_v}\,\kappa^\circ.
\]
Moreover, an ``effective coupling'' form follows by taking the maximum
\emph{before} separating the kernel-side summary statistics:
\[
\widetilde{z}_{\max}
\ \le\
|\mu|\,B_{Y}
\ +\
\sqrt{\widetilde{E}_v}\,\max_{i\neq j}\Big(\sqrt{E_{K}^\circ(i,j)}\,\kappa^{\circ(ij)}\Big).
\]
Defining
\[
\kappa_{\mathrm{eff}}^\circ
\ \defeq\
\max_{i\neq j}\Big(\sqrt{E_{K}^\circ(i,j)}\,\kappa^{\circ(ij)}\Big)
\ =\
\max_{i\neq j}\frac{|\ip{\vX^{\circ(ij)}}{\vY^{(ij)}}|}{\|\vY^{(ij)}\|_2}
\]
we have the deterministic bound
\[
\widetilde{z}_{\max}\ \le\ |\mu|\,B_{Y}\ +\ \sqrt{\widetilde{E}_v}\,\kappa_{\mathrm{eff}}^\circ.
\]
\end{lemma}

\begin{proof}
The decomposition follows by substituting
\(\vX^{(ij)}=\mu\,\vone+\vX^{\circ(ij)}\)
into \(\widetilde{z}_{ij}=\ip{\vX^{(ij)}}{\vY^{(ij)}}\).

For the bound, fix $(i,j)$. By the triangle inequality,
\[
|\widetilde{z}_{ij}|\ \le\ |\mu|\,\big|\ip{\vone}{\vY^{(ij)}}\big|\ +\ \big|\ip{\vX^{\circ(ij)}}{\vY^{(ij)}}\big|.
\]
By definition of the pairwise coupling $\kappa^{\circ(ij)}$,
\(
\big|\ip{\vX^{\circ(ij)}}{\vY^{(ij)}}\big|
\le
\|\vX^{\circ(ij)}\|_{2}\,\|\vY^{(ij)}\|_{2}\,\kappa^{\circ(ij)}.
\)
Since $\|\vY^{(ij)}\|_{2}\le \sqrt{\widetilde{E}_v}$ and $\|\vX^{\circ(ij)}\|_{2}=\sqrt{E_{K}^\circ(i,j)}$, we obtain
\[
|\widetilde{z}_{ij}|
\ \le\
|\mu|\,B_{Y}
\ +\
\sqrt{\widetilde{E}_v}\,\sqrt{E_{K}^\circ(i,j)}\,\kappa^{\circ(ij)}.
\]
Taking $\max_{i\neq j}$ gives
\[
\widetilde{z}_{\max}
\ \le\
|\mu|\,B_{Y}
\ +\
\sqrt{\widetilde{E}_v}\,\max_{i\neq j}\Big(\sqrt{E_{K}^\circ(i,j)}\,\kappa^{\circ(ij)}\Big)
\ =\
|\mu|\,B_{Y}\ +\ \sqrt{\widetilde{E}_v}\,\kappa_{\mathrm{eff}}^\circ.
\]
Finally, since $\sqrt{E_{K}^\circ(i,j)}\le \sqrt{E_{K}^\circ}$ and $\kappa^{\circ(ij)}\le \kappa^\circ$ for all $(i,j)$,
we also have
\(
\kappa_{\mathrm{eff}}^\circ\le \sqrt{E_{K}^\circ}\,\kappa^\circ,
\)
which yields the stated bound
\(
\widetilde{z}_{\max}\le |\mu|\,B_{Y}+\sqrt{E_{K}^\circ}\,\sqrt{\widetilde{E}_v}\,\kappa^\circ.
\)
\end{proof}

\begin{theorem}[Cross-talk bound --- isotropic keys, arbitrary values (bilinear kernel)]
\label{app:regime2}
Assume keys $\vk_{1},\dots,\vk_{F}$ are i.i.d.\ uniform on $\Sph$ and independent of the feature weights
$(\va_{r},\vb_{r})_{r=1}^m$, as defined in the bilinear random features setup (\cref{lem:bilinear-sketched-k2}).
Fix \emph{deterministic} values $\{\vv_{t}\}_{t=1}^F\subset\R^{d}$ and codes $\{\vc_{t}\}_{t=1}^F\subset\R^{d}$.
Fix $\delta\in(0,1)$ and set $L\defeq \log\big(\tfrac{C_{0} F^2}{\delta}\big)$.
Assume $m\le d^2/L$ and let $\mu=1/d$.
Let $\sigma^2\defeq \E\big[(\hat{\matK}_{ti}-\mu)^2\big]=\Theta(\tfrac{1}{d^2}+\tfrac{1}{m})$.
Let $L_v$ be as in \cref{def:Lv}.
Then with probability at least $1-\delta$ (over keys and features),
\begin{equation}
\widetilde{z}_{\max}
\ \le\ \frac{1}{d}\,B_{Y}
\ +\ C\,\sqrt{\widetilde{E}_v}\left(\sigma\sqrt{L}+\frac{L^2}{m}\right)\sqrt{L_v}.
\label{eq:regime2-cross-talk-keff}
\end{equation}
In particular, under $m\le d^2/L$ one has $\sigma\asymp 1/\sqrt{m}$, so
\[
\widetilde{z}_{\max}
\ \le\ \frac{1}{d}\,B_{Y}
\ +\ C\,\sqrt{\widetilde{E}_v}\left(\sqrt{\frac{L}{m}}+\frac{L^2}{m}\right)\sqrt{L_v}.
\]
\end{theorem}

\begin{proof}
We use the effective-coupling reduction from Lemma~\ref{lem:mean-residual-app} and then invoke
Lemma~\ref{lem:keff-iso-keys}.

\paragraph{Deterministic reduction to $\kappa_{\mathrm{eff}}^\circ$.}
Lemma~\ref{lem:mean-residual-app} (with $\mu=1/d$) gives
\[
\widetilde{z}_{\max}
\ \le\
\frac{1}{d}\,B_{Y}
\ +\
\sqrt{\widetilde{E}_v}\,\kappa_{\mathrm{eff}}^\circ.
\]

\paragraph{Concentration of $\kappa_{\mathrm{eff}}^\circ$ (invoke Lemma~\ref{lem:keff-iso-keys}).}
Apply Lemma~\ref{lem:keff-iso-keys} with failure probability $\delta$ to obtain, with probability at least $1-\delta$,
\[
\kappa_{\mathrm{eff}}^\circ
\ \le\
C\left(\sigma\sqrt{L}+\frac{L^2}{m}\right)\sqrt{L_v}.
\]
Substitute this bound into the deterministic reduction to obtain \eqref{eq:regime2-cross-talk-keff}.
\end{proof}


\subsubsection{Isotropic Keys, Isotropic Values (Bilinear Kernel)}


\begin{theorem}[Cross-talk bound --- isotropic keys, isotropic values (bilinear kernel)]
\label{app:isoiso}
Assume \(F\ge3\). Keys \(\vk_1,\dots,\vk_F\) are i.i.d. uniform on \(\Sph\).
Values \(\vv_1,\dots,\vv_F\) are i.i.d. uniform on \(\Sph\) and independent of the
keys and features. The kernel is the bilinear random-feature kernel from
\cref{lem:bilinear-sketched-k2} with \(m\) features.

Fix \(\delta\in(0,1)\), and define
\[
L\defeq \log\left(\frac{C_0F^2}{\delta}\right).
\]
Let
\[
\sigma^2\defeq
\E[(\hat{\matK}_{ti}-1/d)^2]
=
\Theta\left(\frac1{d^2}+\frac1m\right),
\qquad t\neq i.
\]
Assume
\[
m\le \frac{d^2}{L},
\qquad
L^3\le c_0\sigma^2m^2.
\]
Then, with probability at least \(1-\delta\), the following hold simultaneously:
\begin{align*}
B_Y
&\ \le\ C_1\sqrt{\frac{FL}{d}},\\
\widetilde E_v
&\ \le\ C_2\,\frac{(F-2)+L}{d-1},\\
\sqrt{E_K^\circ}
&\ \le\ C_3\,\sigma\sqrt{(F-2)+L},\\
\kappa^\circ
&\ \le\ C_4\sqrt{\frac{L}{F-2}}.
\end{align*}
Consequently,
\[
\widetilde z_{\max}
\le
\frac1dB_Y+
\sqrt{E_K^\circ}\sqrt{\widetilde E_v}\,\kappa^\circ.
\]
In particular, if \(F-2\ge L\), then
\[
\widetilde z_{\max}
\le
C_6\sqrt L\sqrt{\frac{F}{d^3}}
+
C_7\sqrt L\sqrt{\frac{F}{md}}.
\]
\end{theorem}

\begin{proof}
\paragraph{Concentration and union bound.}
Apply Lemma~\ref{lem:BY-iso-values}, Lemma~\ref{lem:Ev-iso-values},
Lemma~\ref{lem:EK-iso-keys}, and Lemma~\ref{lem:kappa-centered-iso-values}, each
with failure probability \(\delta/4\). Increasing \(C_0\) if necessary lets all four
events be written with the same
\[
L=\log\left(\frac{C_0F^2}{\delta}\right).
\]
A union bound gives simultaneous validity with probability at least \(1-\delta\).

\paragraph{Mean--residual reduction.}
On this event, Lemma~\ref{lem:mean-residual-app} with \(\mu=1/d\) gives
\[
\widetilde z_{\max}
\le
\frac1dB_Y+
\sqrt{E_K^\circ}\sqrt{\widetilde E_v}\,\kappa^\circ.
\]

\paragraph{Mean term.}
The mean term satisfies
\[
\frac1dB_Y
\le
C\sqrt{\frac{FL}{d^3}}
=
C\sqrt L\sqrt{\frac{F}{d^3}}.
\]

\paragraph{Residual term.}
Assume now \(F-2\ge L\). By Lemma~\ref{lem:EK-iso-keys},
\[
\sqrt{E_K^\circ}\le C\sigma\sqrt{F-2}.
\]
Therefore
\[
\sqrt{E_K^\circ}\sqrt{\widetilde E_v}\,\kappa^\circ
\le
C\sigma\sqrt{F-2}
\sqrt{\frac{(F-2)+L}{d-1}}
\sqrt{\frac{L}{F-2}}.
\]
Canceling \(\sqrt{F-2}\),
\[
\sqrt{E_K^\circ}\sqrt{\widetilde E_v}\,\kappa^\circ
\le
C\sigma\sqrt L
\sqrt{\frac{(F-2)+L}{d-1}}.
\]
Since \(F-2\ge L\),
\[
(F-2)+L\le2(F-2)\le2F.
\]
Also \(d-1\asymp d\) after adjusting constants. Hence
\[
\sqrt{E_K^\circ}\sqrt{\widetilde E_v}\,\kappa^\circ
\le
C\sigma\sqrt{\frac{FL}{d}}.
\]
Finally, \(m\le d^2/L\) implies
\[
\frac1{d^2}\le \frac1{mL}\le \frac1m,
\]
so
\[
\sigma^2=\Theta\left(\frac1{d^2}+\frac1m\right)
\le
\frac{C}{m}.
\]
Thus
\[
\sigma\sqrt{\frac{FL}{d}}
\le
C\sqrt{\frac{FL}{md}}
=
C\sqrt L\sqrt{\frac{F}{md}}.
\]

\paragraph{Combine.}
Combining the mean and residual estimates gives the final bound.
\end{proof}





\subsection{Signal Bounds}
We now provide lower bounds for the signal term $\widetilde{s}_{ij}$.

\paragraph{Signal as difference of weighted kernel values.}
Recall that under the ``absorbing the competitor'' convention (\eqref{eq:new-signal-crosstalk}), the two-term signal is
\begin{equation}
\widetilde{s}_{ij}
=\hat{\matK}_{ii}\,\ip{\vv_{i}-\vv_j}{\vv_i}
+\hat{\matK}_{ji}\,\ip{\vv_{i}-\vv_j}{\vv_j}
=\hat{\matK}_{ii}\,\ip{\vv_{i}-\vv_j}{\vv_i}
-\hat{\matK}_{ji}\,\ip{\vv_{j}-\vv_i}{\vv_j}.
\label{eq:sij-two-term}
\end{equation}
Since $\hat{\matK}$ is a Gram matrix, we always have
\(
\hat{\matK}_{ii}=\|\phi(\vk_{i})\|_{2}^2\ge 0.
\)
Further, define the smallest (worst-case) signal by
\begin{equation}
\widetilde{s}_{\min}\ \defeq\ \min_{i\neq j}\widetilde{s}_{ij}.
\label{eq:smin}
\end{equation}

\paragraph{Signal summary statistics.}
Define the value-side inner products
\begin{equation}
V_{ij}\ \defeq\ \ip{\vv_{i}-\vv_j}{\vv_i},
\qquad
B_{ij}\ \defeq\ \ip{\vv_{i}-\vv_j}{\vv_j},
\qquad (i\neq j).
\label{eq:AijBij}
\end{equation}
Across our bounds below we will assume, unless stated otherwise,
\begin{equation}
V_{ij}\ \ge\ 0
\qquad\text{for all }i\neq j.
\label{eq:Apos}
\end{equation}
This holds automatically when all values share a common norm:
if $\|\vv_{i}\|_{2}=\|\vv_{j}\|_{2}=r$ then by Cauchy--Schwarz
\(
\ip{\vv_i}{\vv_j}\le \|\vv_{i}\|_{2}\|\vv_{j}\|_{2}=r^2,
\)
hence
\(
V_{ij}=r^2-\ip{\vv_i}{\vv_j}\ge 0.
\)

Recall from \cref{def:Kdiag,def:Vmin} the kernel and value-side extrema:
\begin{equation}
K_{\min}^{\mathrm{diag}}\ =\ \min_{i\in[F]} \hat{\matK}_{ii},
\qquad
K_{\max}^{\mathrm{off}}\ =\ \max_{i\neq j}|\hat{\matK}_{ij}|,
\label{eq:KminKmax}
\end{equation}
\begin{equation}
V_{\min}\ =\ \min_{i\neq j}\ip{\vv_{i}-\vv_j}{\vv_i},
\qquad
V_{\max}\ =\ \max_{i\neq j}\big|\ip{\vv_{i}-\vv_j}{\vv_j}\big|.
\label{eq:AminBmax}
\end{equation}

\subsubsection{Arbitrary Keys, Arbitrary Values}
\label{app:signal-det-arb-arb}

\begin{theorem}[Deterministic signal lower bound --- arbitrary keys, arbitrary values]
Under the positivity condition (\eqref{eq:Apos}), and with the summary statistics (\cref{def:Kdiag,def:Vmin}),
\[
\widetilde{s}_{\min}\ \ge\ K_{\min}^{\mathrm{diag}}\,V_{\min}\ -\ K_{\max}^{\mathrm{off}}\,V_{\max}.
\]
\end{theorem}

\begin{proof}
Fix a pair $(i,j)$ with $i\neq j$.
From \eqref{eq:sij-two-term} and \eqref{eq:AijBij},
\[
\widetilde{s}_{ij}=\hat{\matK}_{ii}\,V_{ij}+\hat{\matK}_{ji}\,B_{ij}.
\]

By definition, $\hat{\matK}_{ii}\ge K_{\min}^{\mathrm{diag}}$.
Also $V_{ij}\ge V_{\min}$ by definition of $V_{\min}$.
Under the positivity assumption (\eqref{eq:Apos}), we have $V_{ij}\ge 0$ for all $i\neq j$, and therefore
multiplying the inequalities $\hat{\matK}_{ii}\ge K_{\min}^{\mathrm{diag}}$ and $V_{ij}\ge V_{\min}$ preserves order:
\[
\hat{\matK}_{ii}\,V_{ij}\ \ge\ K_{\min}^{\mathrm{diag}}\,V_{\min}.
\]

For the second term, we use $xy\ge -|x|\,|y|$:
\[
\hat{\matK}_{ji}\,B_{ij}\ \ge\ -|\hat{\matK}_{ji}|\,|B_{ij}|.
\]
By definition $|\hat{\matK}_{ji}|\le K_{\max}^{\mathrm{off}}$ and $|B_{ij}|\le V_{\max}$, hence
\[
\hat{\matK}_{ji}\,B_{ij}\ \ge\ -K_{\max}^{\mathrm{off}}\,V_{\max}.
\]

Adding the two bounds yields
\(
\widetilde{s}_{ij}\ge K_{\min}^{\mathrm{diag}}V_{\min}-K_{\max}^{\mathrm{off}}V_{\max},
\)
resulting in the claimed bound for $\widetilde{s}_{\min}$.
\end{proof}


\subsubsection{Arbitrary Keys, Isotropic Values}
\begin{theorem}[Signal lower bound --- arbitrary keys, isotropic values]
\label{app:signal-arbkeys-isovalues}
Assume $\vv_{1},\dots,\vv_{F}$ are i.i.d.\ uniform on $\Sph\subset\R^d$, and treat the kernel matrix $\hat{\matK}$ as arbitrary (or condition on it).
Fix $\delta\in(0,1)$ and let
\[
L\ \defeq\ \log\Big(\frac{C_{0} F^2}{\delta}\Big),
\qquad
\varepsilon_{v}\ \defeq\ C_{1}\sqrt{\frac{L}{d-1}}.
\]
Then with probability at least $1-\delta$ (over the values),
\[
\widetilde{s}_{\min}
\ \ge\
K_{\min}^{\mathrm{diag}}\,(1-\varepsilon_{v})
\ -\
K_{\max}^{\mathrm{off}}\,(1+\varepsilon_{v}),
\]
where $K_{\min}^{\mathrm{diag}}$ and $K_{\max}^{\mathrm{off}}$ are the kernel extrema (\cref{def:Kdiag}).
\end{theorem}

\begin{proof}
The proof relies on a single auxiliary concentration bound.

\paragraph{Concentration of the value-side extrema.}
Apply Lemma~\ref{lem:max-pairwise-ip-values} with failure probability $\delta$.
This yields an event $\mathcal{E}_{v}$ such that
$\Pbb(\mathcal{E}_{v}^c)\le \delta$ and on $\mathcal{E}_{v}$,
\[
\max_{i\neq j}\big|\ip{\vv_i}{\vv_j}\big|\ \le\ \varepsilon_{v}.
\]
Hence, on $\mathcal{E}_{v}$,
\[
V_{\min}\ \ge\ 1-\varepsilon_{v},
\qquad
V_{\max}\ \le\ 1+\varepsilon_{v}.
\]

\paragraph{Plug-in.}
On $\mathcal{E}_{v}$, the deterministic signal lower bound gives
\[
\widetilde{s}_{\min}\ \ge\ K_{\min}^{\mathrm{diag}}V_{\min}-K_{\max}^{\mathrm{off}}V_{\max}
\ \ge\ K_{\min}^{\mathrm{diag}}(1-\varepsilon_{v})-K_{\max}^{\mathrm{off}}(1+\varepsilon_{v}),
\]
which is exactly the claimed bound.
\end{proof}


\subsubsection{Isotropic Keys, Arbitrary Values (Bilinear Kernel)}
\begin{theorem}[Signal lower bound --- isotropic keys, arbitrary values (bilinear kernel)]
\label{app:signal-isokeys-arbvalues}
Assume keys $\vk_{1},\dots,\vk_{F}$ are i.i.d.\ uniform on $\Sph$, the kernel is the bilinear random-feature kernel from \cref{lem:bilinear-sketched-k2} with $m$ features, and the values $\{\vv_{t}\}_{t=1}^F$ are deterministic and satisfy the positivity condition (\eqref{eq:Apos}).
Let $V_{\min}$ and $V_{\max}$ be as in \cref{def:Vmin}.
Fix $\delta\in(0,1)$ and set $L\defeq \log(\tfrac{C_{0} F^2}{\delta})$.
Let $\mu\defeq 1/d$ and let $\sigma^2\defeq \E[(\hat{\matK}_{ij}-\mu)^2]=\Theta(\tfrac{1}{d^2}+\tfrac{1}{m})$
(see Lemma~\ref{lem:bilinear-mean-var}).
Then with probability at least $1-\delta$ (over keys and features),
\[
\widetilde{s}_{\min}
\ \ge\
\Big(1-C_{1}\sqrt{\tfrac{L}{m}}\Big)\,V_{\min}
\ -\
\Big(\mu+C_{2}(\sigma\sqrt{L}+L^2/m)\Big)\,V_{\max}.
\]
Moreover, if $L\ge 1$ and $m \geq L^3$,
then the off-diagonal coefficient simplifies and
\[
\widetilde{s}_{\min}
\ \ge\
\Big(1-C_{1}\sqrt{\tfrac{L}{m}}\Big)\,V_{\min}
\ -\
C_{3}\,\sigma\sqrt{L}\,V_{\max}.
\]
\end{theorem}

\begin{proof} The proof combines two auxiliary concentration bounds via a union bound.
\paragraph{Concentration of the diagonal term.}
Apply Lemma~\ref{lem:Kdiag-lower-bilinear} with failure probability $\delta/2$.
This yields an event $\mathcal{E}_{1}$ such that $\Pbb(\mathcal{E}_{1}^c)\le \delta/2$ and on $\mathcal{E}_{1}$,
\[
K_{\min}^{\mathrm{diag}}\ \ge\ 1-C\sqrt{\frac{L}{m}},
\]
after adjusting the absolute constant hidden in $L=\log(\tfrac{C_{0} F^2}{\delta})$.

\paragraph{Concentration of the off-diagonal term.}
Apply Lemma~\ref{lem:Koff-iso-keys} with failure probability $\delta/2$.
This yields an event $\mathcal{E}_{2}$ such that $\Pbb(\mathcal{E}_{2}^c)\le \delta/2$ and on $\mathcal{E}_{2}$,
\[
K_{\max}^{\mathrm{off}}\ \le\ \mu+C\Big(\sigma\sqrt{L}+\frac{L^2}{m}\Big),
\]
again after adjusting the absolute constant hidden in $L$.

\paragraph{Union bound and plug-in.}
By a union bound,
\[
\Pbb(\mathcal{E}_{1}\cap\mathcal{E}_{2})\ \ge\ 1-\Pbb(\mathcal{E}_{1}^c)-\Pbb(\mathcal{E}_{2}^c)\ \ge\ 1-\delta.
\]
On $\mathcal{E}_{1}\cap\mathcal{E}_{2}$, the deterministic signal lower bound gives
\[
\widetilde{s}_{\min}\ \ge\ K_{\min}^{\mathrm{diag}}V_{\min}-K_{\max}^{\mathrm{off}}V_{\max}.
\]
Substituting the bounds for $K_{\min}^{\mathrm{diag}}$ and $K_{\max}^{\mathrm{off}}$ gives, after renaming absolute constants,
\[
\widetilde{s}_{\min}
\ \ge\
\Big(1-C_{1}\sqrt{\tfrac{L}{m}}\Big)V_{\min}
\ -\
\Big(\mu+C_{2}(\sigma\sqrt{L}+L^2/m)\Big)V_{\max},
\]
which is exactly the stated lower bound.

\paragraph{Simplification of the off-diagonal coefficient.}
If \(L\ge 1\) and $m \geq L^3$, then
\(
\sigma^2=\Theta(\tfrac{1}{d^2}+\tfrac{1}{m})
\)
implies
\(
\sigma\gtrsim 1/d=\mu,
\)
so
\(
\mu\le C\sigma\le C\sigma\sqrt{L}.
\)
Hence
\[
\mu+C\Big(\sigma\sqrt{L}+\frac{L^2}{m}\Big)\ \le\ C\sigma\sqrt{L},
\]
and substituting this into the previous bound yields the refined inequality.
\end{proof}


\subsubsection{Isotropic Keys, Isotropic Values (Bilinear Kernel)}
\begin{theorem}[Signal lower bound --- isotropic keys, isotropic values (bilinear kernel)]
\label{app:signal-isokeys-isovalues}
Assume keys $\vk_{1},\dots,\vk_{F}$ are i.i.d.\ uniform on $\Sph$, values $\vv_{1},\dots,\vv_{F}$ are i.i.d.\ uniform on $\Sph$ and independent of the keys and features, and the kernel is the bilinear random-feature kernel from \cref{lem:bilinear-sketched-k2} with $m$ features.
Fix $\delta\in(0,1)$ and set $L\defeq \log(\tfrac{C_{0} F^2}{\delta})$.
Assume $m\le d^2/L$.
Let $\mu\defeq 1/d$ and $\sigma^2\defeq \E[(\hat{\matK}_{ij}-\mu)^2]=\Theta(\tfrac{1}{d^2}+\tfrac{1}{m})$.
Then with probability at least $1-\delta$ (over keys, features, and values),
\[
\widetilde{s}_{\min}
\ \ge\
\Big(1-C_{1}\sqrt{\tfrac{L}{m}}\Big)\,(1-\varepsilon_{v})
\ -\
\Big(\mu+C_{2}(\sigma\sqrt{L}+L^2/m)\Big)\,(1+\varepsilon_{v}),
\qquad
\varepsilon_{v}=C_{3}\sqrt{\tfrac{L}{d-1}}.
\]
Moreover, if \(L\ge 1\) and $m \geq L^3$, then
\[
\widetilde{s}_{\min}
\ \ge\
\Big(1-C_{1}\sqrt{\tfrac{L}{m}}\Big)\,(1-\varepsilon_{v})
\ -\
C_{4}\,\sigma\sqrt{L}\,(1+\varepsilon_{v}),
\qquad
\varepsilon_{v}=C_{3}\sqrt{\tfrac{L}{d-1}}.
\]
\end{theorem}

\begin{proof}
The proof combines three auxiliary concentration bounds via a union bound.
\paragraph{Concentration of the value-side extrema.}
Apply Lemma~\ref{lem:max-pairwise-ip-values} with failure probability $\delta/3$.
This yields an event $\mathcal{E}_{v}$ such that $\Pbb(\mathcal{E}_{v}^c)\le \delta/3$ and on $\mathcal{E}_{v}$,
\[
\max_{i\neq j}\big|\ip{\vv_i}{\vv_j}\big|\ \le\ \varepsilon_{v}.
\]
Hence, on $\mathcal{E}_{v}$,
\[
V_{\min}\ \ge\ 1-\varepsilon_{v},
\qquad
V_{\max}\ \le\ 1+\varepsilon_{v}.
\]

\paragraph{Concentration of the diagonal term.}
Apply Lemma~\ref{lem:Kdiag-lower-bilinear} with failure probability $\delta/3$.
This yields an event $\mathcal{E}_{1}$ such that $\Pbb(\mathcal{E}_{1}^c)\le \delta/3$ and on $\mathcal{E}_{1}$,
\[
K_{\min}^{\mathrm{diag}}\ \ge\ 1-C\sqrt{\frac{L}{m}},
\]
after adjusting the absolute constant hidden in $L$.

\paragraph{Concentration of the off-diagonal term.}
Apply Lemma~\ref{lem:Koff-iso-keys} with failure probability $\delta/3$.
This yields an event $\mathcal{E}_{2}$ such that $\Pbb(\mathcal{E}_{2}^c)\le \delta/3$ and on $\mathcal{E}_{2}$,
\[
K_{\max}^{\mathrm{off}}\ \le\ \mu+C\Big(\sigma\sqrt{L}+\frac{L^2}{m}\Big),
\]
again after adjusting the absolute constant hidden in $L$.

\paragraph{Union bound and plug-in.}
By a union bound,
\[
\Pbb(\mathcal{E}_{v}\cap\mathcal{E}_{1}\cap\mathcal{E}_{2})
\ \ge\ 1-\Pbb(\mathcal{E}_{v}^c)-\Pbb(\mathcal{E}_{1}^c)-\Pbb(\mathcal{E}_{2}^c)
\ \ge\ 1-\delta.
\]
On this intersection, the deterministic signal lower bound gives
\[
\widetilde{s}_{\min}\ \ge\ K_{\min}^{\mathrm{diag}}V_{\min}-K_{\max}^{\mathrm{off}}V_{\max}.
\]
Substituting the bounds for $K_{\min}^{\mathrm{diag}}$, $K_{\max}^{\mathrm{off}}$, $V_{\min}$, and $V_{\max}$ yields the stated lower bound.

\paragraph{Simplification of the off-diagonal coefficient.}
If \(L\ge 1\) and $m \geq L^3$, then the same reasoning as above gives
\[
\mu+C\Big(\sigma\sqrt{L}+\frac{L^2}{m}\Big)\ \le\ C\sigma\sqrt{L},
\]
and substituting this estimate into the previous bound yields the refined inequality.
\end{proof}

\subsection{Margin Bounds}
We now combine our signal and cross-talk bounds to bound the Hebbian memory margin.
\subsubsection{Arbitrary Keys, Arbitrary Values}
\label{app:combined-deterministic}
\label{app:theory:margin_bounds:akav}
\phantomsection\label{app:theory:proof:arb-arb-det}
\begin{theorem}[Margin bound --- arbitrary keys, arbitrary values]
\label{thm:combined-deterministic}
Assume the positivity condition~\eqref{eq:Apos} (i.e., \(V_{ij}\ge 0\) for all \(i\neq j\)).
Let \(K_{\min}^{\mathrm{diag}}\), \(K_{\max}^{\mathrm{off}}\) (\cref{def:Kdiag}), \(V_{\min}\), and \(V_{\max}\) (\cref{def:Vmin}) be the signal-side summary statistics, and let \(\widetilde{E}_{K}\), \(\widetilde{E}_{v}\), and \(\widetilde{\kappa}\) be the
cross-talk summary statistics from \eqref{eq:summary-stats}.
Then
\begin{equation}
\gamma_{\min}\ \ge\ K_{\min}^{\mathrm{diag}}\,V_{\min}\ -\ K_{\max}^{\mathrm{off}}\,V_{\max}\ -\ \sqrt{\widetilde{E}_K}\,\sqrt{\widetilde{E}_v}\,\widetilde{\kappa}.
\label{eq:combined-det-bound}
\end{equation}
Moreover, if we instead keep the original signal/cross-talk decomposition (\eqref{eq:signal-crosstalk})
(i.e., without absorbing the competitor term), then applying Cauchy--Schwarz with the coupling factor
$\kappa$ from \cref{def:kappa} yields the alternative bound
\begin{equation}
\gamma_{\min}\ \ge\ K_{\min}^{\mathrm{diag}}\,V_{\min}\ -\ \sqrt{E_K}\,\sqrt{E_v}\,\kappa,
\label{eq:app:akav:thm}
\end{equation}
where \(E_{K}\) and \(E_{v}\) are as defined in \cref{def:Ecol,def:EY}.

\end{theorem}
\begin{proof}
Under the positivity condition (\eqref{eq:Apos}), the deterministic signal lower bound gives
\[
\widetilde{s}_{\min}\ \ge\ K_{\min}^{\mathrm{diag}}\,V_{\min}\ -\ K_{\max}^{\mathrm{off}}\,V_{\max}.
\]

Further, Theorem~\ref{lem:deterministic} yields
\[
\widetilde{z}_{\max}\ \le\ \sqrt{\widetilde{E}_K}\,\sqrt{\widetilde{E}_v}\,\widetilde{\kappa}.
\]

Recalling that \(\gamma_{\min}\ge \widetilde{s}_{\min}-\widetilde{z}_{\max}\), we therefore have
\[
\gamma_{\min}
\ \ge\ \widetilde{s}_{\min}-\widetilde{z}_{\max}
\ \ge\ K_{\min}^{\mathrm{diag}}\,V_{\min}
-
K_{\max}^{\mathrm{off}}\,V_{\max}
-
\sqrt{\widetilde{E}_K}\,\sqrt{\widetilde{E}_v}\,\widetilde{\kappa},
\]


Finally, to obtain \eqref{eq:app:akav:thm}, start from the original decomposition \eqref{eq:signal-crosstalk}:
\[
\gamma_{ij} \,=\, \hat{\matK}_{ii}\,V_{ij} \,+\, z_{ij},
\qquad
z_{ij}\ \defeq\ \sum_{t\neq i}\hat{\matK}_{ti}\,\ip{\vv_{i}-\vv_j}{\vc_t}.
\]
By definition of the coupling factor $\kappa^{ij}$ (\cref{def:kappa}),
\(
|z_{ij}|\ \le\ \sqrt{E_{K}(i)}\,\sqrt{E_{v}(i,j)}\,\kappa^{ij}
\ \le\ \sqrt{E_K}\,\sqrt{E_v}\,\kappa.
\)
Since $\hat{\matK}_{ii}\ge K_{\min}^{\mathrm{diag}}$ and $V_{ij}\ge V_{\min}$, taking $\min_{i\neq j}$ yields \eqref{eq:app:akav:thm}.

\end{proof}

\begin{corollary}[Arbitrary-embedding margin scaling --- simplified]\label{cor:bilinear-scaling-simplified}
Under the assumptions of \Cref{thm:combined-deterministic}, further reasonably assuming $d\ge 4\log F$ and using the penalty statistics $S_{\mathrm{sig}}$, $P_{\mathrm{key}}$, $P_{\mathrm{val}}$, $P_{\mathrm{align}}$ from \cref{sec:arbk-isov}, for some $C\in(0,1)$ we have
\[
\gamma_{\min}
\;\geq\;
C\,S_{\mathrm{sig}}
\;-\;
P_{\mathrm{key}}\,P_{\mathrm{val}}\,P_{\mathrm{align}}\sqrt{\frac{F\log F}{md}}.
\]
\end{corollary}
\begin{proof}
The statistics $E_K,E_v,\kappa$ (\cref{sec:arbk-isov}) sum over $t\neq i$ (the competitor $t=j$ included), so they coincide with the untruncated statistics of \cref{def:Ecol,def:EY,def:kappa} and \Cref{thm:combined-deterministic} (\cref{eq:app:akav:thm}) applies:
\(
\gamma_{\min}\ge K_{\min}^{\mathrm{diag}}\,V_{\min}-\sqrt{E_K}\,\sqrt{E_v}\,\kappa.
\)
Substituting $\sqrt{E_K}=P_{\mathrm{key}}\sqrt{F/m}$, $\sqrt{E_v}=P_{\mathrm{val}}\sqrt{F/d}$, and $\kappa=P_{\mathrm{align}}\sqrt{\log(F)/F}$ gives the exact identity
\[
\sqrt{E_K}\,\sqrt{E_v}\,\kappa
= P_{\mathrm{key}}\,P_{\mathrm{val}}\,P_{\mathrm{align}}\sqrt{\tfrac{F}{m}\cdot\tfrac{F}{d}\cdot\tfrac{\log F}{F}}
= P_{\mathrm{key}}\,P_{\mathrm{val}}\,P_{\mathrm{align}}\sqrt{\tfrac{F\log F}{md}}.
\]
For the signal, the definition of $S_{\mathrm{sig}}$ gives $K_{\min}^{\mathrm{diag}}\,V_{\min}=\big(1-\sqrt{\log(F)/d}\big)\,S_{\mathrm{sig}}$. Since $d\ge 4\log F$ forces $\sqrt{\log(F)/d}\le\tfrac12$, and $S_{\mathrm{sig}}\ge 0$ (as $K_{\min}^{\mathrm{diag}}\ge 0$ and $V_{\min}\ge 0$ by~\cref{eq:Apos}), we get $K_{\min}^{\mathrm{diag}}\,V_{\min}\ge C\,S_{\mathrm{sig}}$ with $C\defeq\tfrac12$. Combining the two equations proves the claim.
\end{proof}

\begin{corollary}[Arbitrary-embedding fact-storage capacity]\label{cor:optimal-fact-storage-capacity-arb-formal}
Under the assumptions of \Cref{cor:bilinear-scaling-simplified}, the margin lower bound is positive precisely when
\[
P_{\mathrm{key}}\,P_{\mathrm{val}}\,P_{\mathrm{align}}\sqrt{\frac{F\log F}{md}} < C\,S_{\mathrm{sig}},
\qquad\text{i.e.}\qquad
md > \frac{1}{C^2}\,F\log F\left(\frac{P_{\mathrm{key}} P_{\mathrm{val}} P_{\mathrm{align}}}{S_{\mathrm{sig}}}\right)^{\!2}.
\]
Hence our construction stores all $F$ facts with positive margin using a parameter budget
\[
W = md = \Theta\!\left(F\log F\left(\frac{P_{\mathrm{key}} P_{\mathrm{val}} P_{\mathrm{align}}}{S_{\mathrm{sig}}}\right)^{\!2}\right),
\]
the same rate fact-storage rate as in the isotropic setting up to the penalization factors.
\end{corollary}


\subsubsection{Arbitrary Keys, Isotropic Values}
\label{app:combined-regime1}
\label{app:theory:margin_bounds:akrv}
\phantomsection\label{app:theory:proof:arbk-isov}

\begin{theorem}[Margin bound --- arbitrary keys, isotropic values]\label{thm:combined-akrv}
Assume \(\vv_{1},\dots,\vv_{F}\) are i.i.d. uniform on \(\Sph\subset\R^d\), and treat the kernel matrix \(\hat{\matK}\) as arbitrary.
Fix \(\delta\in(0,1)\) and let
\[
L\ \defeq\ \log\!\Big(\frac{C_{0} F^2}{\delta}\Big),
\qquad
\varepsilon_{v}\ \defeq\ C_{1}\sqrt{\frac{L}{d-1}}.
\]
Define the (untruncated) kernel-column energy
\[
E_{K}\ \defeq\ \max_{i\in[F]}\ \sum_{t\neq i}\hat{\matK}_{ti}^2.
\]
Then with probability at least \(1-\delta\) (over the values),
\begin{equation}
\gamma_{\min}\ \ge\ K_{\min}^{\mathrm{diag}}(1-\varepsilon_{v})-K_{\max}^{\mathrm{off}} - C_{2}\sqrt{E_K}\,\sqrt{\frac{L}{d-1}}\,\sqrt{1+\frac{L}{F-2}}.
\label{eq:app:akrv:thm-full}
\end{equation}
In particular, if \(F-2\ge L\), then
\begin{equation}
\gamma_{\min}
\ \ge\
K_{\min}^{\mathrm{diag}}(1-\varepsilon_{v})-K_{\max}^{\mathrm{off}}
\ -\ C_{3}\sqrt{E_K}\,\sqrt{\frac{L}{d}}.
\label{eq:app:akrv:thm}
\end{equation}
\end{theorem}

\begin{proof}
The proof combines the signal and cross-talk bounds via a union bound.
\paragraph{Signal lower bound.}
Apply the signal bound from Theorem~\ref{app:signal-arbkeys-isovalues} with failure probability \(\delta/2\).
This yields an event \(\mathcal E_{\mathrm{sig}}\) such that
\(
\Pbb(\mathcal E_{\mathrm{sig}}^c)\le \delta/2
\)
and on \(\mathcal E_{\mathrm{sig}}\),
\begin{equation}
\widetilde{s}_{\min}
\ \ge\
K_{\min}^{\mathrm{diag}}(1-\varepsilon_{v})
\ -\
K_{\max}^{\mathrm{off}}(1+\varepsilon_{v}),
\label{eq:combined-reg1-signal}
\end{equation}
after the standard adjustment of the absolute constant hidden in
\(
L=\log(\tfrac{C_{0} F^2}{\delta})
\)
so that replacing \(\delta\) by \(\delta/2\) only changes the constant.

\paragraph{Cross-talk bound.}
Apply the cross-talk bound from Theorem~\ref{app:regime1} with failure probability \(\delta/2\).
This yields an event \(\mathcal E_{\mathrm{cross}}\) such that
\(
\Pbb(\mathcal E_{\mathrm{cross}}^c)\le \delta/2
\)
and on \(\mathcal E_{\mathrm{cross}}\),
\begin{equation}
\widetilde{z}_{\max}
\ \le\
C\sqrt{E_K}\,\sqrt{\frac{L}{d-1}}\,\sqrt{1+\frac{L}{F-2}}.
\label{eq:combined-reg1-xtalk}
\end{equation}
Here we used the trivial bound \(\widetilde{E}_{K}\le E_{K}\), since dropping coordinates can only decrease the $\ell_{2}$ norm.
Again we have absorbed the harmless \(\delta\mapsto \delta/2\) change into the absolute constant inside \(L\).

\paragraph{Union bound and combine.}
By a union bound,
\[
\Pbb(\mathcal E_{\mathrm{sig}}\cap \mathcal E_{\mathrm{cross}})
\ \ge\
1-\Pbb(\mathcal E_{\mathrm{sig}}^c)-\Pbb(\mathcal E_{\mathrm{cross}}^c)
\ \ge\ 1-\delta.
\]
On \(\mathcal E_{\mathrm{sig}}\cap \mathcal E_{\mathrm{cross}}\), use \(\gamma_{\min}\ge \widetilde{s}_{\min}-\widetilde{z}_{\max}\).
Substituting \eqref{eq:combined-reg1-signal} and \eqref{eq:combined-reg1-xtalk} yields
\begin{align}
\gamma_{\min}
&\ge
K_{\min}^{\mathrm{diag}}(1-\varepsilon_{v})
-
K_{\max}^{\mathrm{off}}(1+\varepsilon_{v})
-
C_{1}\sqrt{E_K}\,\sqrt{\frac{L}{d-1}}\,\sqrt{1+\frac{L}{F-2}}
\notag\\
&=
K_{\min}^{\mathrm{diag}}(1-\varepsilon_{v})
-
K_{\max}^{\mathrm{off}}
-
\Bigg[
K_{\max}^{\mathrm{off}}\,\varepsilon_{v}
+
C_{1}\sqrt{E_K}\,\sqrt{\frac{L}{d-1}}\,\sqrt{1+\frac{L}{F-2}}
\Bigg].
\label{eq:combined-reg1-before-absorb}
\end{align}

\paragraph{Absorb the competitor term.}
We claim that
\begin{equation}
K_{\max}^{\mathrm{off}}\ \le\ \sqrt{E_K}.
\label{eq:Koff-le-sqrtEK}
\end{equation}
Choose indices \((a,b)\) with \(a\neq b\) such that
\(
|\hat{\matK}_{ba}|=K_{\max}^{\mathrm{off}}.
\)
Then by definition of \(E_{K}\),
\[
E_{K}
\ \ge\
\sum_{t\neq a}\hat{\matK}_{ta}^2
\ \ge\
\hat{\matK}_{ba}^2
\ =\
\big(K_{\max}^{\mathrm{off}}\big)^2,
\]
which proves \eqref{eq:Koff-le-sqrtEK}.
Therefore,
\[
K_{\max}^{\mathrm{off}}\,\varepsilon_{v}
\ \le\
\sqrt{E_K}\,\varepsilon_{v}
\ =\
C_{2}\sqrt{E_K}\,\sqrt{\frac{L}{d-1}}
\ \le\
C_{3}\sqrt{E_K}\,\sqrt{\frac{L}{d-1}}\,\sqrt{1+\frac{L}{F-2}},
\]
since \(\sqrt{1+L/(F-2)}\ge 1\).
Absorbing this into the bracket in \eqref{eq:combined-reg1-before-absorb} yields the claimed bound.

\paragraph{Simplified form when \(F-2\ge L\).}
If \(F-2\ge L\), then \(\sqrt{1+L/(F-2)}\le \sqrt{2}\).
Also \(d-1\asymp d\) for \(d\ge 2\), so after adjusting the absolute constant we obtain the simplified bound stated in the theorem.
\end{proof}


\subsubsection{Isotropic Keys, Arbitrary Values (Bilinear Kernel)}
\label{app:combined-regime2}
\label{app:theory:margin_bounds:rkav}
\phantomsection\label{app:theory:proof:isok-bilinear-hp}

\begin{theorem}[Margin bound --- isotropic keys, arbitrary values (bilinear kernel)]\label{thm:combined-rkav}
Assume keys \(\vk_{1},\dots,\vk_{F}\) are i.i.d. uniform on \(\Sph\), the kernel is the bilinear random-feature kernel from \cref{lem:bilinear-sketched-k2} with $m$ features, the values \(\vv_{1},\dots,\vv_{F}\in\R^d\) are deterministic, and the positivity condition \(\ip{\vv_{i}-\vv_j}{\vv_i}\ge 0\) holds for all \(i\neq j\).
Let
\[
\mu\ \defeq\ \frac{1}{d},
\qquad
\sigma^2\ \defeq\ \E[(\hat{\matK}_{ti}-\mu)^2]=\Theta\!\Big(\frac{1}{d^2}+\frac{1}{m}\Big),
\]
and fix \(\delta\in(0,1)\). Set
\[
L\ \defeq\ \log\!\Big(\frac{C_{0} F^2}{\delta}\Big), \qquad \varepsilon_k \defeq \sqrt{\frac{L}{m}}
\]
and assume in addition that \(F-2\ge L\), \(m \geq L^3\), and \(m\le d^2/L\).

Recall the centered-sum budget \(B_Y\) (\cref{def:BY}) and let \(L_v\) be as in \cref{def:Lv}.
Then with probability at least \(1-\delta\) (over the keys and features),
\begin{equation}
\gamma_{\min}
\ \ge\
\Big(1-C_{1}\varepsilon_k\Big)V_{\min}
-
C_{2}\Big(B_{Y}+\sqrt{E_v}\,\sqrt{L_v}\Big)
\Big(\sigma\sqrt{L}\Big).
\label{eq:app:rkav:thm}
\end{equation}
\end{theorem}

\begin{proof}
The proof combines the signal and cross-talk bounds via a union bound.
\paragraph{Signal lower bound.}
Apply the signal bound from Theorem~\ref{app:signal-isokeys-arbvalues} with failure probability \(\delta/2\).
This yields an event \(\mathcal E_{\mathrm{sig}}\) such that
\(
\Pbb(\mathcal E_{\mathrm{sig}}^c)\le \delta/2
\)
and on \(\mathcal E_{\mathrm{sig}}\),
\begin{equation}
\widetilde{s}_{\min}
\ \ge\
\Big(1-C_{1}\sqrt{\tfrac{L}{m}}\Big)V_{\min}
-
\Big(\frac{1}{d}+C_{2}\big(\sigma\sqrt{L}+\tfrac{L^2}{m}\big)\Big)V_{\max}.
\label{eq:combined-reg2-signal}
\end{equation}
As before, replacing \(\delta\) by \(\delta/2\) only changes the absolute constant hidden in \(L\).

\paragraph{Cross-talk bound.}
Apply the cross-talk bound from Theorem~\ref{app:regime2} with failure probability \(\delta/2\).
This yields an event \(\mathcal E_{\mathrm{cross}}\) such that
\(
\Pbb(\mathcal E_{\mathrm{cross}}^c)\le \delta/2
\)
and on \(\mathcal E_{\mathrm{cross}}\),
\begin{equation}
\widetilde{z}_{\max}
\ \le\
\frac{1}{d}B_{Y}
\ +\ C_{1}\sqrt{\widetilde{E}_v}\Big(\sigma\sqrt{L}+\frac{L^2}{m}\Big)\sqrt{L_v}.
\label{eq:combined-reg2-xtalk}
\end{equation}
After adjusting the absolute constant,
\begin{equation}
\widetilde{z}_{\max}
\ \le\
\frac{1}{d}B_{Y}
\ +\ C_{2}\sqrt{\widetilde{E}_v}\Big(\sigma\sqrt{L}+\frac{L^2}{m}\Big)\sqrt{L_v}.
\label{eq:combined-reg2-xtalk-simple}
\end{equation}

\paragraph{Union bound and combine.}
By a union bound,
\[
\Pbb(\mathcal E_{\mathrm{sig}}\cap\mathcal E_{\mathrm{cross}})\ \ge\ 1-\delta.
\]
On this intersection, use \(\gamma_{\min}\ge \widetilde{s}_{\min}-\widetilde{z}_{\max}\).
Substituting \eqref{eq:combined-reg2-signal} and \eqref{eq:combined-reg2-xtalk-simple} gives
\begin{align}
\gamma_{\min}
&\ge
\Big(1-C_{1}\sqrt{\tfrac{L}{m}}\Big)V_{\min}
-
\frac{1}{d}(B_{Y}+V_{\max})
\notag\\
&\hspace{1.6em}
-
C_{2}\Big(V_{\max}+\sqrt{\widetilde{E}_v}\,\sqrt{L_v}\Big)\Big(\sigma\sqrt{L}+\frac{L^2}{m}\Big).
\label{eq:combined-reg2-combined}
\end{align}

\paragraph{Absorb the value-side terms.}
Since \(V_{\max}\le \sqrt{E_v}\) and \(\sqrt{\widetilde{E}_v}\le \sqrt{E_v}\), and \(L_{v}\ge1\) by \cref{def:Lv} (which holds even in the degenerate case \(\widetilde{E}_{v}=0\), where the cross-talk term vanishes but the signal-side \(V_{\max}\) need not, by the convention \(L_v=1\)), we have \(V_{\max}+\sqrt{\widetilde{E}_v}\sqrt{L_v}\le \sqrt{E_v}\sqrt{L_v}+\sqrt{E_v}\sqrt{L_v}= 2\sqrt{E_v}\sqrt{L_v}\). Therefore, after adjusting the absolute constant,
\[
C_{2}\Big(V_{\max}+\sqrt{\widetilde{E}_v}\,\sqrt{L_v}\Big)\Big(\sigma\sqrt{L}+\frac{L^2}{m}\Big)
\ \le\
C_{3}\sqrt{E_v}\,\sqrt{L_v}\Big(\sigma\sqrt{L}+\frac{L^2}{m}\Big).
\]


\paragraph{Absorb the \(d^{-1}\) factor.}
Since \(L=\log(\tfrac{C_{0} F^2}{\delta})\ge 1\) and
\(
\sigma^2=\Theta\!\big(\tfrac{1}{d^2}+\tfrac{1}{m}\big)
\)
implies \(\sigma\gtrsim 1/d\), we have \(\frac{1}{d}\le C\,\sigma\sqrt{L}\).
Therefore, using also \(V_{\max}\le \sqrt{E_v}\le \sqrt{E_v}\,\sqrt{L_v}\),
\[
\frac{1}{d}(B_{Y}+V_{\max})
\ \le\
C_{1}\Big(B_{Y}+\sqrt{E_v}\,\sqrt{L_v}\Big)\sigma\sqrt{L}
\ \le\
C_{2}\Big(B_{Y}+\sqrt{E_v}\,\sqrt{L_v}\Big)\Big(\sigma\sqrt{L}+\frac{L^2}{m}\Big).
\]
Plugging the last two estimates into \eqref{eq:combined-reg2-combined} and using $m \geq L^3$, we have
\(\gamma_{\min} \ge \Big(1-C_{4}\sqrt{\tfrac{L}{m}}\Big)V_{\min} - C_{5}\Big(B_{Y}+\sqrt{E_v}\,\sqrt{L_v}\Big)\sigma\sqrt{L}\),
which is exactly the stated bound.
\end{proof}


\subsubsection{Isotropic Keys, Isotropic Values (Bilinear Kernel)}
\label{app:combined-regime3}
\label{app:theory:margin_bounds:rkrv}
\phantomsection\label{app:theory:proof:exactK2}
\phantomsection\label{app:theory:proof:bilinear-scaling}
\phantomsection\label{app:theory:proof:bilinear-fact-capacity}

\begin{theorem}[Margin bound --- isotropic keys, isotropic values (bilinear kernel)]\label{thm:combined-rkrv}
Assume keys \(\vk_{1},\dots,\vk_{F}\) are i.i.d. uniform on \(\Sph\), values \(\vv_{1},\dots,\vv_{F}\) are i.i.d. uniform on \(\Sph\) and independent of the keys, and the kernel is the bilinear random-feature kernel from \cref{lem:bilinear-sketched-k2} with $m$ features.
Fix \(\delta\in(0,1)\) and set
\[
L\ \defeq\ \log\!\Big(\frac{C_{0} F^2}{\delta}\Big).
\]
Assume $F \ge 2L$, $m \ge L^3$, and $m\le d^2/L$. Then, with probability at least \(1-\delta\) (over keys, features, and values), the minimum ``absorbed-competitor'' margin satisfies
\begin{equation}
\gamma_{\min}
\ \ge\
1
\;-\;
C_{1}\sqrt{\frac{L}{m}}
\;-\;
\varepsilon_{v}
\;-\;
C_{3}\,\sqrt{L}\sqrt{\frac{F}{d^3}}
\;-\;
C_{4}\,\sqrt{L}\sqrt{\frac{F}{md}},
\qquad
\varepsilon_{v}\defeq C_{2}\sqrt{\frac{L}{d-1}}.
\label{eq:app:rkrv:sketchedk2_full}
\end{equation}
\end{theorem}


\begin{proof}
\paragraph{Signal lower bound.}
Apply the signal bound from Theorem~\ref{app:signal-isokeys-isovalues} with failure probability $\delta/2$. Since $m\ge L^3$, this yields an event $\mathcal E_{\mathrm{sig}}$ with $\Pbb(\mathcal E_{\mathrm{sig}}^c)\le\delta/2$ on which
\[
\widetilde{s}_{\min}
\ \ge\
\Big(1-C_{1}\sqrt{\frac{L}{m}}\Big)(1-\varepsilon_{v})
\ -\
C_{4}\,\sigma\sqrt{L}\,(1+\varepsilon_{v}),
\]
with $\varepsilon_{v}$ as in the theorem statement, after the standard adjustment of the absolute constant hidden in $L$.

\paragraph{Cross-talk bound.}
Apply the cross-talk bound from Theorem~\ref{app:isoiso} with failure probability $\delta/2$. Since $F\ge 2L$ gives $F-2\ge L$, this yields an event $\mathcal E_{\mathrm{cross}}$ with $\Pbb(\mathcal E_{\mathrm{cross}}^c)\le\delta/2$ on which
\[
\widetilde{z}_{\max}
\ \le\
C_{6}\,\sqrt L\sqrt{\frac{F}{d^3}}
\ +\
C_{7}\,\sqrt L\sqrt{\frac{F}{md}}.
\]

\paragraph{Union bound and combine.}
By a union bound, $\Pbb(\mathcal E_{\mathrm{sig}}\cap\mathcal E_{\mathrm{cross}})\ge 1-\delta$. On this intersection, $\gamma_{\min}\ge\widetilde{s}_{\min}-\widetilde{z}_{\max}$ gives
\[
\gamma_{\min}
\ \ge\
\Big(1-C_{1}\sqrt{\frac{L}{m}}\Big)(1-\varepsilon_{v})
\ -\
C_{4}\,\sigma\sqrt{L}\,(1+\varepsilon_{v})
\ -\
C_{6}\,\sqrt L\sqrt{\frac{F}{d^3}}
\ -\
C_{7}\,\sqrt L\sqrt{\frac{F}{md}}.
\]
Since $m\le d^2/L$ forces $\sigma\asymp 1/\sqrt m$, and $\varepsilon_{v}=O(1)$, the signal off-diagonal penalty obeys $C_{4}\,\sigma\sqrt{L}\,(1+\varepsilon_{v})\le C\sqrt{L/m}$ and is absorbed into the first term. Expanding $\big(1-C\sqrt{L/m}\big)(1-\varepsilon_{v})\ge 1-C\sqrt{L/m}-\varepsilon_{v}$ and renaming the absolute constants yields \cref{eq:app:rkrv:sketchedk2_full}.
\end{proof}

\begin{corollary}[Isotropic margin scaling --- simplified]\label{cor:informal-iso-iso-margin}
Under the assumptions of \Cref{thm:combined-rkrv} and $m,d\le F$, retaining only the dominant cross-talk term gives, with probability at least $1-\delta$,
\[
\gamma_{\min}
\;\geq\;
1
\;-\;
C\sqrt{\frac{F\log(F/\delta)}{md}}.
\]
\end{corollary}
\begin{proof}
The dominant cross-talk term in \cref{eq:app:rkrv:sketchedk2_full} is $C_{4}\sqrt{L}\sqrt{F/(md)}$; under the standing assumptions, the remaining terms are lower order and are absorbed into the constant $C$. Since $L = \log(C_0 F^2/\delta) = \Theta(\log(F/\delta))$, we have $\sqrt{L}\sqrt{F/(md)} = \Theta\!\big(\sqrt{F\log(F/\delta)/(md)}\big)$, which yields the stated form.
\end{proof}

\begin{corollary}[Isotropic fact-storage capacity]\label{cor:optimal-fact-storage-capacity-formal}
Under the assumptions of \Cref{cor:informal-iso-iso-margin}, the margin lower bound is positive precisely when
\[
C\sqrt{\frac{F\log(F/\delta)}{md}} < 1,
\qquad\text{i.e.}\qquad
md > C^2\,F\log(F/\delta).
\]
Hence our construction stores all $F$ facts with positive margin using a parameter budget
\[
W = md = \Theta\!\left(F\log F\right),
\]
the information-theoretically optimal fact-storage rate.
\end{corollary}

\subsubsection{Summary Table: Margin Bounds Across Regimes}

\begin{table}[H]
\centering
\renewcommand{\arraystretch}{1.5}
\begin{tabular}{c|c|c}
&
\makecell{\textbf{Random}\ \textbf{keys}}
&
\makecell{\textbf{Arbitrary}\ \textbf{keys}}
\\ \hline
\makecell{\textbf{Random}\\ \textbf{values}}
&
$\begin{aligned}
\gamma_{\min} &\ge\; 1 \\
&-\, C\sqrt{\frac{FL}{md}}
\end{aligned}$
&
$\begin{aligned}
\gamma_{\min} &\ge\; K_{\min}^{\mathrm{diag}}(1-\varepsilon_{v})-K_{\max}^{\mathrm{off}} \\
&-\, C\sqrt{E_K}\,\sqrt{\frac{L}{d}}
\end{aligned}$
\\ \hline
\makecell{\textbf{Arbitrary}\\ \textbf{values}}
&
$\begin{aligned}
\gamma_{\min} &\ge\; \Big(1-\varepsilon_{k}\Big)V_{\min} \\
&-\, C\!\left(B_Y+\sqrt{E_v}\sqrt{L_v}\right)\!\sqrt{\frac{L}{m}}
\end{aligned}$
&
$\begin{aligned}
\gamma_{\min} &\ge\; K_{\min}^{\mathrm{diag}}V_{\min} \\
&-\, \sqrt{E_K}\,\sqrt{E_v}\,\kappa
\end{aligned}$
\end{tabular}
\caption{\textbf{Margin bounds across key/value geometry regimes.}
Each cell shows the decoding margin scaling for a given key/value geometry.
The top-left (isotropic) cell is the baseline.
Relaxing the key geometry introduces key-geometry statistics ($K_{\min}^{\mathrm{diag}}$, $K_{\max}^{\mathrm{off}}$, $E_{K}$);
relaxing value geometry introduces value-geometry statistics ($V_{\min}$, $B_{Y}$, $V_{\max}$, $E_{v}$, $L_{v}$);
relaxing both geometries introduces a coupling factor $\kappa$.
The arbitrary keys and values cell uses the raw summary-statistic form of \Cref{thm:combined-deterministic}, which is equivalent to the penalty-statistic general margin/capacity bound in \Cref{cor:optimal-fact-storage-capacity-arb-formal} of the main text.
See Theorems~\ref{thm:combined-deterministic}, \ref{thm:combined-akrv}, \ref{thm:combined-rkav} and \ref{thm:combined-rkrv} for formal statements and Appendix~\ref{app:exp:section4:margins} for empirical scaling results.}
\label{tab:2x2}
\end{table}


\subsubsection{A Welch/Frobenius Upper Bound in the Isotropic Key/Value Regime}
\phantomsection\label{app:theory:margin_bounds:rkrv:welch}
\label{app:welch-upper-bound-iso}

In the main text, \Cref{thm:bilinear-scaling} provides a lower bound on the minimum margin achieved by our construction in the isotropic key/value regime.
We now show that this lower bound is asymptotically tight, up to constants and logarithmic factors, by proving a matching \emph{upper bound} on the best possible margin of any admissible rank-limited kernel memory in this regime. The key idea is that the arbitrary-keys / isotropic-values analysis in \Cref{app:theory:margin_bounds:akrv} is controlled by the column energy
\[
E_{\mathrm{col}}(i) := \sum_{t\neq i} \hat{\matK}_{ti}^2.
\]
Thus, to upper bound the achievable margin, it suffices to show that for any rank-$m$ kernel, some column must have nontrivial off-diagonal squared mass.
We achieve this using a Welch / Frobenius argument: intuitively, low rank forces a non-negligible amount of kernel mass off the diagonal, which in turn yields an unavoidable cross-talk floor.

We begin by formalizing the class of kernels under consideration.

\begin{definition}[Rank-$m$ kernel]
\label{def:rank-m-kernel}
Let $\vk_{1},\dots,\vk_{F} \in \mathbb{R}^d$ be the stored keys. We say that
$\hat{\matK} \in \mathbb{R}^{F\times F}$ is a \emph{rank-$m$ kernel} if there exists a feature
map $\phi(\cdot)$ with $\phi(\vk_{i})\in\mathbb{R}^m$ such that
\[
\hat{\matK}_{ij} = \langle \phi(\vk_{i}), \phi(\vk_{j})\rangle
\qquad\text{for all } i,j\in[F].
\]
Equivalently, $\hat{\matK} \succeq 0$ and $\operatorname{rank}(\hat{\matK})\le m$.
\end{definition}

The following lemma is the kernel-side ingredient. It lower bounds the maximum column energy
of any PSD rank-$m$ kernel in terms of its diagonal.

\begin{lemma}[Welch/Frobenius lower bound for rank-$m$ kernels]
\label{lem:welch-frobenius-general}
Let $\hat{\matK} \in \mathbb{R}^{F\times F}$ be a rank-$m$ kernel in the sense of
Definition~\ref{def:rank-m-kernel}. Define
\[
E_{\mathrm{col}}(i) := \sum_{t\neq i}\hat{\matK}_{ti}^2,
\qquad
E_{\mathrm{col}} := \max_{i\in[F]} E_{\mathrm{col}}(i).
\]
Then
\[
\|\hat{\matK}\|_{F}^2 \ge \frac{\operatorname{tr}(\hat{\matK})^2}{m},
\]
and consequently
\[
E_{\mathrm{col}}
\;\ge\;
\frac{1}{F}\left(
\frac{\operatorname{tr}(\hat{\matK})^2}{m}
-
\sum_{i=1}^F \hat{\matK}_{ii}^2
\right).
\]
In particular, if
\[
\hat{\matK}_{ii}\in[1-\varepsilon_{\mathrm{diag}},\,1+\varepsilon_{\mathrm{diag}}]
\qquad\text{for all } i\in[F],
\]
then
\[
E_{\mathrm{col}}
\;\ge\;
\frac{F(1-\varepsilon_{\mathrm{diag}})^2}{m}
-
(1+\varepsilon_{\mathrm{diag}})^2.
\]
\end{lemma}

\begin{proof}
Since $\hat{\matK} \succeq 0$ and $\operatorname{rank}(\hat{\matK})\le m$, let
$\lambda_{1},\dots,\lambda_{r}$ be the nonzero eigenvalues of $\hat{\matK}$, where $r\le m$. Then
\[
\operatorname{tr}(\hat{\matK})=\sum_{a=1}^r \lambda_{a},
\qquad
\|\hat{\matK}\|_{F}^2=\sum_{a=1}^r \lambda_{a}^2.
\]
Applying Cauchy--Schwarz to $(\lambda_{1},\dots,\lambda_{r})$ gives
\[
\Big(\sum_{a=1}^r \lambda_{a}\Big)^2
\le r\sum_{a=1}^r \lambda_{a}^2
\le m\sum_{a=1}^r \lambda_{a}^2,
\]
hence
\[
\|\hat{\matK}\|_{F}^2 \ge \frac{\operatorname{tr}(\hat{\matK})^2}{m}.
\]

Next expand the Frobenius norm entrywise:
\[
\|\hat{\matK}\|_{F}^2
=
\sum_{i=1}^F \hat{\matK}_{ii}^2
+
\sum_{i=1}^F\sum_{t\neq i}\hat{\matK}_{ti}^2
=
\sum_{i=1}^F \hat{\matK}_{ii}^2
+
\sum_{i=1}^F E_{\mathrm{col}}(i).
\]
Therefore
\[
\sum_{i=1}^F E_{\mathrm{col}}(i)
\ge
\frac{\operatorname{tr}(\hat{\matK})^2}{m}
-
\sum_{i=1}^F \hat{\matK}_{ii}^2.
\]
Since the maximum is at least the average,
\[
E_{\mathrm{col}}
=
\max_{i} E_{\mathrm{col}}(i)
\ge
\frac{1}{F}\sum_{i=1}^F E_{\mathrm{col}}(i)
\ge
\frac{1}{F}\left(
\frac{\operatorname{tr}(\hat{\matK})^2}{m}
-
\sum_{i=1}^F \hat{\matK}_{ii}^2
\right).
\]

If moreover $\hat{\matK}_{ii}\in[1-\varepsilon_{\mathrm{diag}},1+\varepsilon_{\mathrm{diag}}]$ for all $i$,
then
\[
\operatorname{tr}(\hat{\matK})=\sum_{i=1}^F \hat{\matK}_{ii} \ge F(1-\varepsilon_{\mathrm{diag}})
\]
and
\[
\sum_{i=1}^F \hat{\matK}_{ii}^2 \le F(1+\varepsilon_{\mathrm{diag}})^2.
\]
Substituting these into the previous display yields
\[
E_{\mathrm{col}}
\ge
\frac{1}{F}\left(
\frac{F^2(1-\varepsilon_{\mathrm{diag}})^2}{m}
-
F(1+\varepsilon_{\mathrm{diag}})^2
\right)
=
\frac{F(1-\varepsilon_{\mathrm{diag}})^2}{m}
-
(1+\varepsilon_{\mathrm{diag}})^2.
\]
\end{proof}

Regarding the assumption
\[
\hat{\matK}_{ii}\in[1-\varepsilon_{\mathrm{diag}},1+\varepsilon_{\mathrm{diag}}],
\]
note that for the exact quadratic kernel on unit-norm isotropic keys, one has
\[
K_{2}(\vk_{i},\vk_{i})=\langle \vk_{i},\vk_{i}\rangle^2 = 1
\]
exactly.
For the sketched / bilinear random-feature kernel used in our construction, the diagonal is generally not exactly $1$, but it instead concentrates near $1$ with high probability by \Cref{lem:Kdiag-lower-bilinear}.

Combining the kernel-side lower bound with the isotropic-value competitor argument used in the
proof of \Cref{rem:welch-benchmark} yields the desired asymptotic optimality statement.

\begin{corollary}[Rank-$m$ kernels are asymptotically unbeatable in the isotropic key/value regime]
\label{cor:rank-m-unbeatable-iso}
Assume the isotropic key/value setting, and let $\hat{\matK}$ be any rank-$m$ kernel independent of
the values, with
\[
\hat{\matK}_{ii}\in[1-\varepsilon_{\mathrm{diag}},\,1+\varepsilon_{\mathrm{diag}}]
\qquad\text{for all } i\in[F].
\]
Then the isotropic-value competitor argument used in the proof of \Cref{rem:welch-benchmark} yields
\[
\gamma_{\min}
\;\le\;
(1+\varepsilon_{\mathrm{diag}})
-
\Omega\!\left(
\sqrt{\frac{E_{\mathrm{col}}\log F}{d}}
\right)
+
\mathcal{O}\!\left(
\sqrt{\frac{\log F}{d}}
\right)
+
\mathcal{O}\!\left(
\sqrt{\frac{E_{\mathrm{col}}\log(1/\delta)}{d}}
\right)
+
\mathcal{O}\!\left(
\sqrt{\frac{E_{\mathrm{col}}}{d}}
\right).
\]
Combining with Lemma~\ref{lem:welch-frobenius-general} yields
\[
\gamma_{\min}
\;\le\;
(1+\varepsilon_{\mathrm{diag}})
-
\Omega\!\left(
\sqrt{
\frac{\big(\frac{F(1-\varepsilon_{\mathrm{diag}})^2}{m}-(1+\varepsilon_{\mathrm{diag}})^2\big)\log F}{d}
}
\right)
+ \text{lower-order terms}.
\]
In particular, if $\varepsilon_{\mathrm{diag}}=o(1)$ and $F/m\to\infty$, then up to constants and
logarithmic factors,
\[
\gamma_{\min}
\;\le\;
1 - \Omega\!\left(\sqrt{\frac{(F/m)\log F}{d}}\right).
\]
Equivalently, no admissible rank-$m$ kernel can asymptotically beat the
$\sqrt{(F/m)\log F / d}$ margin floor in the isotropic key/value regime, up to constants and logs.
\end{corollary}

\begin{proof}
Apply the isotropic-value competitor argument used in the proof of \Cref{rem:welch-benchmark} to
the column $i^\star$ with $E_{\mathrm{col}}(i^\star)=E_{\mathrm{col}}$. This gives
\[
\gamma_{\min}
\le
\hat{\matK}_{i^\star i^\star}
-
\Omega\!\left(
\sqrt{\frac{E_{\mathrm{col}}\log F}{d}}
\right)
+
\mathcal{O}\!\left(
|\hat{\matK}_{i^\star i^\star}|\sqrt{\frac{\log F}{d}}
\right)
+
\mathcal{O}\!\left(
\sqrt{\frac{E_{\mathrm{col}}\log(1/\delta)}{d}}
\right)
+
\mathcal{O}\!\left(
\sqrt{\frac{E_{\mathrm{col}}}{d}}
\right),
\]
after absorbing harmless constant-factor differences into the big-$\mathcal{O}$/big-$\Omega$
notation.

Using $\hat{\matK}_{i^\star i^\star}\le 1+\varepsilon_{\mathrm{diag}}$ gives the first bound.
The second bound results from substituting the lower bound on $E_{\mathrm{col}}$ from~\Cref{lem:welch-frobenius-general}.

Finally, if $\varepsilon_{\mathrm{diag}}=o(1)$, then
\[
\frac{F(1-\varepsilon_{\mathrm{diag}})^2}{m}
-
(1+\varepsilon_{\mathrm{diag}})^2
=
\Theta\!\left(\frac{F}{m}\right)
\]
whenever $F/m\to\infty$, so the dominant negative term scales as
\[
\sqrt{\frac{(F/m)\log F}{d}},
\]
which proves the last claim.
\end{proof}


\subsection{Auxiliary Results}


\subsubsection{Signal Bounds: Isotropic Values}
\label{app:max-pairwise-ip-values}

\begin{lemma}[Isotropic inner-product concentration]
\label{lem:iso-innerprod}
Let $\vv\sim\mathrm{Unif}(\Sph)$ with $d\ge 2$.
Then for any fixed $\vw\in\R^d$ and all $t\ge 0$,
$\|\ip{\vw}{\vv}\|_{\psi_2}\lesssim \|\vw\|_{2}/\sqrt{d-1}$ and
$\big\|\ip{\vw}{\vv}^2-\E[\ip{\vw}{\vv}^2]\big\|_{\psi_1}\ \lesssim\ \frac{\norm{\vw}_{2}^2}{d-1}$. I.e.\ $\ip{\vw}{\vv}$ is sub-Gaussian and $\ip{\vw}{\vv}^2-\E[\ip{\vw}{\vv}^2]$ is sub-exponential, with the corresponding scale parameters.
\end{lemma}
\begin{proof}
    This is a standard consequence of concentration of measure on the sphere, see \cite{vershynin2026high}.
\end{proof}


\begin{lemma}[Max pairwise inner product for isotropic values]
\label{lem:max-pairwise-ip-values}
Assume $\vv_{1},\dots,\vv_{F}$ are i.i.d.\ uniform on $\Sph$ with $d\ge 2$.
Fix $\delta\in(0,1)$ and set
\begin{equation}
L\ \defeq\ \log\Big(\frac{C_{0} F^2}{\delta}\Big)
\label{eq:L-signal-values}
\end{equation}
for a sufficiently large absolute constant $C_{0}$.
Then with probability at least $1-\delta$,
\begin{equation}
\max_{i\neq j}\big|\ip{\vv_i}{\vv_j}\big|
\ \le\
C_{1}\,\sqrt{\frac{L}{d-1}}.
\label{eq:maxpair-ip-values}
\end{equation}
\end{lemma}

\begin{proof}
Fix a pair $(i,j)$ with $i\neq j$.
Condition on $\vv_{i}$ and define $\mathcal{F}_{i}\defeq\sigma(\vv_{i})$.
Then $\vv_{j}$ is still uniform on $\Sph$ and independent of $\mathcal{F}_{i}$.
Applying Lemma~\ref{lem:iso-innerprod} with $\vw=\vv_{i}$ (note $\|\vv_{i}\|_{2}=1$) yields that conditional on $\mathcal{F}_{i}$,
the scalar $\ip{\vv_i}{\vv_j}$ is centered sub-Gaussian with scale $\lesssim 1/\sqrt{d-1}$.
Equivalently, there exist absolute constants $c,C_{0}>0$ such that for all $t\ge 0$,
\begin{equation}
\Pbb\Big(\big|\ip{\vv_i}{\vv_j}\big|\ge t\ \Big|\ \mathcal{F}_{i}\Big)
\ \le\ 2e^{-c(d-1)t^2}.
\label{eq:pair-ip-tail}
\end{equation}
Since the bound from \eqref{eq:pair-ip-tail} does not depend on the realized $\vv_{i}$,
it also holds unconditionally:
\[
\Pbb\Big(\big|\ip{\vv_i}{\vv_j}\big|\ge t\Big)\ \le\ 2e^{-c(d-1)t^2}.
\]

Now take a union bound over all pairs $(i,j)$ with $i\neq j$ (at most $F^2$ pairs):
\[
\Pbb\Big(\max_{i\neq j}\big|\ip{\vv_i}{\vv_j}\big|\ge t\Big)
\ \le\
\sum_{i\neq j}\Pbb\Big(\big|\ip{\vv_i}{\vv_j}\big|\ge t\Big)
\ \le\
2F^2 e^{-c(d-1)t^2}.
\]
Choose
\(
t\defeq C\sqrt{L/(d-1)}
\)
with $C$ large enough so that $2F^2 e^{(-c(d-1)t^2)}\le\delta$.
This yields \eqref{eq:maxpair-ip-values}.
\end{proof}




\subsubsection{Signal Bounds: Isotropic Keys and Bilinear Random Features}
\label{app:Koff-iso-keys}

\begin{lemma}[Uniform off-diagonal bound (isotropic keys, bilinear kernel)]
\label{lem:Koff-iso-keys}
Assume keys are i.i.d.\ uniform on $\Sph$ and the kernel is the bilinear random-feature kernel from \cref{lem:bilinear-sketched-k2} with $m$ features.
Let $\mu\defeq\E[\hat{\matK}_{ij}]=1/d$ for $i\neq j$ and let $\sigma^2\defeq\E[(\hat{\matK}_{ij}-\mu)^2]$ as in
Lemma~\ref{lem:bilinear-mean-var}.
Fix $\delta\in(0,1)$ and define $L\defeq \log(\tfrac{C_{0} F^2}{\delta})$.
Assume that $m \le d^2/L$.
Then with probability at least $1-\delta$ (over keys and features),
\begin{equation}
\max_{i\neq j}\big|\hat{\matK}_{ij}-\mu\big|
\ \le\
C_{1}\Big(\sigma\sqrt{L}+\frac{L^2}{m}\Big),
\label{eq:Koff-centered-max}
\end{equation}
and hence
\begin{equation}
K_{\max}^{\mathrm{off}}=\max_{i\neq j}|\hat{\matK}_{ij}|
\ \le\
\mu\ +\ C_{2}\Big(\sigma\sqrt{L}+\frac{L^2}{m}\Big).
\label{eq:Koff-abs-max}
\end{equation}
\end{lemma}

\begin{proof}
For any fixed pair $(i,j)$ with $i\neq j$, apply Lemma~\ref{lem:chaos-centered-mu} with failure probability $\delta' \coloneqq \delta/F^2$.
With $L' \coloneqq \log(6/\delta') = \log(6F^2/\delta)\le L$ (for $C_{0}$ large enough), the assumption $m\le d^2/L$ implies $m \le d^2/L'$, so Lemma~\ref{lem:chaos-centered-mu} yields
\[
    \big|\hat{\matK}_{ij}-\mu\big|
    \le
    C\Big(\sqrt{\tfrac{L'}{m}} + \tfrac{(L')^2}{m} + \tfrac{\sqrt{L'}}{d}\Big)
\]
with probability at least $1-\delta'$.
A union bound over the at most $F^2$ off-diagonal pairs gives that, with probability at least $1-\delta$,
\[
    \max_{i\neq j}\big|\hat{\matK}_{ij}-\mu\big|
    \le
    C\Big(\sqrt{\tfrac{L}{m}} + \tfrac{L^2}{m} + \tfrac{\sqrt{L}}{d}\Big).
\]
Finally, since $\sigma^2=\Theta(1/d^2+1/m)$ by Lemma~\ref{lem:bilinear-mean-var}, we have $\sqrt{L/m}+\sqrt{L}/d \lesssim \sigma\sqrt{L}$, proving \eqref{eq:Koff-centered-max}.
Equation~\eqref{eq:Koff-abs-max} follows by the triangle inequality.
\end{proof}


\label{app:Kdiag-lower-bilinear}
\begin{lemma}[Uniform diagonal lower bound (bilinear kernel)]
\label{lem:Kdiag-lower-bilinear}
Assume the bilinear random-feature kernel from \cref{lem:bilinear-sketched-k2} with $m$ features, and assume $\|\vk_{i}\|_{2}=1$.
Then there exist absolute constants $c,C_{0}>0$ such that for every $t\in(0,1)$ and every $i\in[F]$,
\begin{equation}
\Pbb\big(\hat{\matK}_{ii}\le 1-t\big)\ \le\ e^{-c m t^2}.
\label{eq:Kdiag-lower-tail}
\end{equation}
Consequently, for any $\delta\in(0,1)$ and $L\defeq \log(\tfrac{C_{0} F^2}{\delta})$, with probability at least $1-\delta$,
\begin{equation}
K_{\min}^{\mathrm{diag}}=\min_{i\in[F]}\hat{\matK}_{ii}
\ \ge\
1\ -\ C_{1}\sqrt{\frac{L}{m}}.
\label{eq:Kdiag-min}
\end{equation}
\end{lemma}

\begin{proof}
Fix $i$ and write
\[
\hat{\matK}_{ii}=\frac1m\sum_{r=1}^m Z_{r},\qquad Z_{r}:=\big((\va_{r}^\top \vk_{i})(\vb_{r}^\top \vk_{i})\big)^2
=(G_rH_{r})^2,
\]
where $G_{r},H_{r}\stackrel{\text{i.i.d.}}{\sim}\mathcal N(0,1)$. Then
$\E [Z_{r}]=1$ and $\Var(Z_{r})=\E[Z_{r}^2]-(\E[Z_{r}])^2=9-1=8$.
Let $Y_{r}:=1-Z_{r}$. Then $\E[Y_{r}]=0$, $\Var(Y_{r})=8$, and since $Z_{r}\ge0$ we have $Y_{r}\le 1$. Therefore, for $t\in(0,1)$,
\[
\Pbb(\hat{\matK}_{ii}\le 1-t)
=\Pbb\!\left(\sum_{r=1}^m Y_{r}\ge mt\right).
\]
By the one-sided Bernstein inequality for independent mean-zero variables bounded above by $1$,
\[
\Pbb\!\left(\sum_{r=1}^m Y_{r}\ge mt\right)
\le
\exp\!\left(-\frac{(mt)^2}{2(8m+\frac{mt}{3})}\right)
=
\exp\!\left(-\frac{m t^2}{2(8+t/3)}\right)
\le
e^{-cmt^2}
\]
for an absolute $c>0$, proving \eqref{eq:Kdiag-lower-tail}.
The uniform bound (\eqref{eq:Kdiag-min}) follows by a union bound over $i\in[F]$.
\end{proof}

\subsubsection{Cross-Talk Bounds: Isotropic Values}

\begin{lemma}[Mean-term concentration $B_{Y}$ (isotropic values)]
\label{lem:BY-iso-values}
Assume $\vv_{1},\dots,\vv_{F}$ are i.i.d.\ uniform on $\Sph$.
Fix $\delta\in(0,1)$ and let $L\defeq \log(\tfrac{C_{0} F^2}{\delta})$.
Then with probability at least $1-\delta$ (over the values),
\begin{equation}
B_{Y}\ \defeq\ \max_{i\in[F]}\max_{j\neq i}\big|\ip{\vone}{\vY^{(ij)}}\big|
\ \le\ C_{1}\sqrt{\frac{FL}{d}}.
\label{eq:BY-iso-values}
\end{equation}
\end{lemma}

\begin{proof}
Fix a pair $(i,j)$ with $j\neq i$ and set $\vw_{ij}\defeq \vv_{i}-\vv_{j}$.
Let
\[
\mathcal{F}_{ij}\ \defeq\ \sigma(\vv_{i},\vv_{j}).
\]
Conditional on $\mathcal{F}_{ij}$, the random variables
\(
\theta_{t}\defeq \ip{\vw_{ij}}{\vv_t}
\)
for $t\notin\{i,j\}$ are independent, centered, and sub-Gaussian with
$\|\theta_{t}\|_{\psi_2}\lesssim \|\vw_{ij}\|_{2}/\sqrt{d-1}$ by Lemma~\ref{lem:iso-innerprod}.
Therefore, conditional on $\mathcal{F}_{ij}$, the sum
\[
\ip{\vone}{\vY^{(ij)}}=\sum_{t\notin\{i,j\}}\theta_{t}
\]
is centered sub-Gaussian with
\(
\big\|\ip{\vone}{\vY^{(ij)}}\big\|_{\psi_2}\lesssim \|\vw_{ij}\|_{2}\sqrt{\tfrac{F-2}{d-1}}
\lesssim \sqrt{\tfrac{F}{d}}.
\)
Hence for all $s\ge 0$,
\[
\Pbb\Big(
\big|\ip{\vone}{\vY^{(ij)}}\big|\ \ge\ C\sqrt{\tfrac{Fs}{d}}
\ \Big|\ \mathcal{F}_{ij}
\Big)
\ \le\ 2e^{-s}.
\]
This conditional tail bound is uniform in $\mathcal{F}_{ij}$, so taking expectations over $\mathcal{F}_{ij}$
yields the same inequality unconditionally.

Now choose $s= L=\log(\tfrac{C_{0} F^2}{\delta})$ so that
$2e^{-s}\le \delta/F^2$.
Union bound over all pairs $(i,j)$ (at most $F(F-1)\le F^2$ choices) gives \eqref{eq:BY-iso-values}.
\end{proof}


\begin{lemma}[Value-difference energy concentration $\widetilde{E}_{v}$ (isotropic values)]
\label{lem:Ev-iso-values}
Assume $\vv_{1},\dots,\vv_{F}$ are i.i.d.\ uniform on $\Sph$.
Fix $\delta\in(0,1)$ and let $L\defeq \log(\tfrac{C_{0} F^2}{\delta})$.
Then with probability at least $1-\delta$ (over the values),
\begin{equation}
\widetilde{E}_{v}\ \defeq\ \max_{i\in[F]}\max_{j\neq i}\sum_{t\notin\{i,j\}}\ip{\vv_{i}-\vv_j}{\vv_t}^2
\ \le\ C_{1}\,\frac{(F-2)+L}{d-1}.
\label{eq:Ev-iso-values}
\end{equation}
\end{lemma}

\begin{proof}
Fix a pair $(i,j)$ with $j\neq i$ and write $\vw_{ij}\defeq \vv_{i}-\vv_{j}$.
Let
\(
\mathcal{F}_{ij}\defeq \sigma(\vv_{i},\vv_{j})
\)
and, for $t\notin\{i,j\}$, define $\theta_{t}\defeq \ip{\vw_{ij}}{\vv_t}$.
Conditional on $\mathcal{F}_{ij}$, the $\theta_{t}$ are independent, centered, sub-Gaussian, therefore the centered squares
\[
Z_{t}\defeq \theta_{t}^2-\E[\theta_{t}^2\mid \mathcal{F}_{ij}]
\]
are independent and sub-exponential with
\(
\|Z_{t}\|_{\psi_1}\lesssim \|\vw_{ij}\|_{2}^2/(d-1)
\)
by Lemma~\ref{lem:iso-innerprod}.

Write
\[
\widetilde{E}_{v}(i,j)\ \defeq\ \sum_{t\notin\{i,j\}}\theta_{t}^2
\ =\ \sum_{t\notin\{i,j\}}\E[\theta_{t}^2\mid\mathcal{F}_{ij}]
\ +\ \sum_{t\notin\{i,j\}}Z_{t}.
\]
Since $\vv_{t}$ is isotropic, $\E[\theta_{t}^2\mid\mathcal{F}_{ij}]=\|\vw_{ij}\|_{2}^2/d$ and hence
$\sum_{t\notin\{i,j\}}\E[\theta_{t}^2\mid\mathcal{F}_{ij}]=(F-2)\|\vw_{ij}\|_{2}^2/d$.
Also $\|\vw_{ij}\|_{2}\le 2$.
A standard Bernstein inequality for sums of independent sub-exponential variables implies that for all $s\ge 0$,
\[
\Pbb\Big(
\sum_{t\notin\{i,j\}}Z_{t}\ \ge\ C\,\frac{\|\vw_{ij}\|_{2}^2}{d-1}\big(\sqrt{(F-2)s}+s\big)
\ \Big|\ \mathcal{F}_{ij}
\Big)
\ \le\ 2e^{-s}.
\]
As in Lemma~\ref{lem:BY-iso-values}, the bound is uniform in $\mathcal{F}_{ij}$, so the same tail holds unconditionally.

Using $\sqrt{(F-2)s}\le \tfrac{(F-2)+s}{2}$ and absorbing constants yields
\[
\Pbb\Big(
\widetilde{E}_{v}(i,j)\ \ge\ C\,\frac{\|\vw_{ij}\|_{2}^2}{d-1}\big((F-2)+s\big)
\Big)
\ \le\ 2e^{-cs}.
\]
Since $\|\vw_{ij}\|_{2}^2\le 4$, choosing $s= L=\log(\tfrac{C_{0} F^2}{\delta})$ makes the right-hand side
$\le \delta/F^2$.
Union bounding over all pairs $(i,j)$ yields \eqref{eq:Ev-iso-values}.
\end{proof}


\begin{lemma}[Coupling concentration $\widetilde{\kappa}$ (isotropic values)]
\label{lem:kappa-iso-values}
Assume $\vv_{1},\dots,\vv_{F}$ are i.i.d.\ uniform on $\Sph$.
Treat the kernel matrix $\hat{\matK}$ as arbitrary/deterministic.
Fix $\delta\in(0,1)$ and let $L\defeq \log(\tfrac{C_{0} F^2}{\delta})$.
Then with probability at least $1-\delta$ (over the values), the coupling statistic $\widetilde{\kappa}$ from \eqref{eq:summary-stats} satisfies
\begin{equation}
\widetilde{\kappa}\ \le\ C_{1}\,\sqrt{\frac{L}{F-2}}.
\label{eq:kappa-iso-values}
\end{equation}
\end{lemma}

\begin{proof}
Fix a pair $(i,j)$ with $j\neq i$.
If $\|\vX^{(ij)}\|_{2}=0$ or $\|\vY^{(ij)}\|_{2}=0$ then $\widetilde{\kappa}^{ij}=0$, so assume both norms are positive and set
\[
\vu^{(ij)}\defeq \vX^{(ij)}/\|\vX^{(ij)}\|_{2}.
\]
Write $\vw_{ij}\defeq \vv_{i}-\vv_{j}$ and let $\mathcal{F}_{ij}\defeq \sigma(\vv_{i},\vv_{j})$.
For $t\notin\{i,j\}$ define $\theta_{t}\defeq \ip{\vw_{ij}}{\vv_t}$. Let
\[
S_{ij}\ \defeq\ \ip{\vu^{(ij)}}{\vY^{(ij)}}\ =\ \sum_{t\notin\{i,j\}}\vu^{(ij)}_{t}\,\theta_{t}.
\]
Then
\[
\widetilde{\kappa}^{ij}=\frac{|\ip{\vX^{(ij)}}{\vY^{(ij)}}|}{\|\vX^{(ij)}\|_{2}\|\vY^{(ij)}\|_2}
=\frac{|\ip{\vu^{(ij)}}{\vY^{(ij)}}|}{\|\vY^{(ij)}\|_2}
=\frac{|S_{ij}|}{\|\vY^{(ij)}\|_2}.
\]

\paragraph{Numerator tail.}
Conditional on $\mathcal{F}_{ij}$, the vectors $(\vv_{t})_{t\notin\{i,j\}}$ are i.i.d.\ uniform on $\Sph$,
so each $\theta_{t}=\ip{\vw_{ij}}{\vv_t}$ is centered sub-Gaussian with
$\|\theta_{t}\|_{\psi_2}\lesssim \|\vw_{ij}\|_{2}/\sqrt{d-1}$ by Lemma~\ref{lem:iso-innerprod}.
Since $\sum_{t\notin\{i,j\}}(\vu^{(ij)}_{t})^2\le 1$, $S_{ij}$ is centered sub-Gaussian with
$\|S_{ij}\|_{\psi_2}\lesssim \|\vw_{ij}\|_{2}/\sqrt{d-1}$. Hence for all $s\ge 0$,
\[
\Pbb\Big(|S_{ij}|\ \ge\ C\,\|\vw_{ij}\|_{2}\sqrt{\tfrac{s}{d-1}}\ \Big|\ \mathcal{F}_{ij}\Big)\ \le\ 2e^{-s}.
\]
The right-hand side does not depend on $\mathcal{F}_{ij}$, so taking expectation over $\mathcal{F}_{ij}$ yields the same bound unconditionally.

\paragraph{Denominator lower tail.}
Conditional on $\mathcal{F}_{ij}$, the centered squares
$Z_{t}\defeq \theta_{t}^2-\E[\theta_{t}^2\mid\mathcal{F}_{ij}]$ are independent, mean-zero, and sub-exponential with
$\|Z_{t}\|_{\psi_1}\lesssim \|\vw_{ij}\|_{2}^2/(d-1)$ (Lemma~\ref{lem:iso-innerprod}).
A Bernstein inequality for sums of independent sub-exponential variables (applied to $-Z_{t}$) yields that for all $s\ge 0$,
\[
\Pbb\Bigg(
\|\vY^{(ij)}\|_{2}^2
\le
\sum_{t\notin\{i,j\}}\E[\theta_{t}^2\mid\mathcal{F}_{ij}]
-
C\,\frac{\|\vw_{ij}\|_{2}^2}{d-1}\big(\sqrt{(F-2)s}+s\big)
\ \Bigg|\ \mathcal{F}_{ij}
\Bigg)
\ \le\ e^{-s}.
\]
Again the bound is uniform in $\mathcal{F}_{ij}$, so it holds unconditionally after taking the expectation.
Since $\vv_{t}$ is isotropic, $\E[\theta_{t}^2\mid\mathcal{F}_{ij}]=\|\vw_{ij}\|_{2}^2/d$, and hence
$\sum_{t\notin\{i,j\}}\E[\theta_{t}^2\mid\mathcal{F}_{ij}]=(F-2)\|\vw_{ij}\|_{2}^2/d$.
Using $1/d\ge 1/(2(d-1))$ for $d\ge 2$, the complement event implies
\[
\|\vY^{(ij)}\|_{2}^2
\ge
\frac{\|\vw_{ij}\|_{2}^2}{d-1}\Big(\frac{F-2}{2}-C\big(\sqrt{(F-2)s}+s\big)\Big).
\]

\paragraph{A high-probability ratio bound.}
On the intersection of the numerator and denominator events, if $F-2\ge C_{0} s$,
then $\|\vY^{(ij)}\|_{2}\ge c\,\|\vw_{ij}\|_{2}\sqrt{(F-2)/(d-1)}$ and hence
\[
\widetilde{\kappa}^{ij}=\frac{|S_{ij}|}{\|\vY^{(ij)}\|_2}\ \le\ C\,\sqrt{\frac{s}{F-2}}.
\]
If instead $F-2< C_{0} s$, then $\sqrt{s/(F-2)}\ge c$ and the same bound holds trivially since $\widetilde{\kappa}^{ij}\le 1$.
Overall, for the fixed $(i,j)$ the bound holds with failure probability at most $3e^{-s}$.

\paragraph{Choose $s$ and union bound.}
Set $s=L=\log(\tfrac{C_{0} F^2}{\delta})$.
Then for a fixed $(i,j)$ the failure probability is $\le 3e^{-L}\le \delta/F^2$ (for $C_{0}$ large enough).
Union bounding over all pairs $(i,j)$ yields \eqref{eq:kappa-iso-values}.
\end{proof}


\begin{lemma}[Centered coupling concentration \(\kappa^\circ\) (isotropic values)]
\label{lem:kappa-centered-iso-values}
Assume \(\vv_1,\dots,\vv_F\) are i.i.d. uniform on \(\Sph\). Let
\(\vX^{\circ(ij)}\in\R^{F-2}\), \(i\neq j\), be deterministic vectors or random
vectors independent of the values. Define
\[
\kappa^{\circ(ij)}
\defeq
\frac{\left|\left\langle \vX^{\circ(ij)},\vY^{(ij)}\right\rangle\right|}
{\|\vX^{\circ(ij)}\|_2\|\vY^{(ij)}\|_2},
\qquad
\kappa^\circ\defeq\max_{i\neq j}\kappa^{\circ(ij)},
\]
with the convention \(\kappa^{\circ(ij)}=0\) if a denominator vanishes. Fix
\(\delta\in(0,1)\), and set
\[
L\defeq \log\left(\frac{C_0F^2}{\delta}\right).
\]
Then, with probability at least \(1-\delta\),
\[
\kappa^\circ\le C\sqrt{\frac{L}{F-2}}.
\]
\end{lemma}

\begin{proof}
Condition on \(\{\vX^{\circ(ij)}:i\neq j\}\). After conditioning, these columns
are deterministic and independent of the values. The proof of
Lemma~\ref{lem:kappa-iso-values} applies unchanged with \(\vX^{(ij)}\) replaced
by \(\vX^{\circ(ij)}\). Removing the conditioning proves the claim.
\end{proof}


\subsubsection{Cross-Talk Bounds: Isotropic Keys and Bilinear Random Features}
\begin{lemma}[Bilinear kernel mean and variance (isotropic keys)]
\label{lem:bilinear-mean-var}
Let $\hat{\matK}$ be the bilinear random-feature kernel from \cref{lem:bilinear-sketched-k2} with $m$ features and
Gaussian weights $(\va_{r},\vb_{r})_{r=1}^m$.
Assume keys $\vk_{1},\dots,\vk_{F}$ are i.i.d.\ uniform on $\Sph$.
Fix $t\neq i$ and write $\rho=\ip{\vk_t}{\vk_i}$.
Then:
\begin{enumerate}[leftmargin=2.2em]
\item (Conditional mean) $\E[\hat{\matK}_{ti}\mid \rho]=\rho^2$.
\item (Unconditional mean) $\mu\defeq \E[\hat{\matK}_{ti}]=\E[\rho^2]=1/d$.
\item (Variance scale) letting $\sigma^2\defeq \E[(\hat{\matK}_{ti}-\mu)^2]$, we have
$\sigma^2=\Theta(\tfrac{1}{d^2}+\tfrac{1}{m})$ %
\end{enumerate}
\end{lemma}

\begin{proof}

Condition on $(\vk_{t},\vk_{i})$ and hence on $\rho$.
Write
\[
\hat{\matK}_{ti}=\frac{1}{m}\sum_{r=1}^m U_{r},
\qquad
U_{r} \defeq (\va_{r}^\top \vk_{t})(\va_{r}^\top \vk_{i})\,(\vb_{r}^\top \vk_{t})(\vb_{r}^\top \vk_{i}),
\]
with $(\va_{r},\vb_{r})$ i.i.d.\ standard Gaussian.
Let $A_{r}\defeq (\va_{r}^\top \vk_{t})(\va_{r}^\top \vk_{i})$ and $B_{r}\defeq (\vb_{r}^\top \vk_{t})(\vb_{r}^\top \vk_{i})$ so that $U_{r}=A_rB_{r}$
and $A_{r}$ is independent of $B_{r}$.

\emph{Mean.}
Since $(\va_{r}^\top \vk_{t},\va_{r}^\top \vk_{i})$ is centered bivariate Gaussian with covariance $\rho$,
$\E[A_{r}\mid \rho]=\rho$ and likewise $\E[B_{r}\mid\rho]=\rho$, hence
$\E[U_{r}\mid\rho]=\rho^2$ and $\E[\hat{\matK}_{ti}\mid\rho]=\rho^2$.
Averaging over isotropic keys yields $\mu=\E[\rho^2]=1/d$.

\emph{Variance scale.}
A Wick/Isserlis calculation gives $\E[A_{r}^2\mid\rho]=1+2\rho^2$ and hence
$\Var(U_{r}\mid\rho)=1+4\rho^2+3\rho^4$.
Therefore
$\Var(\hat{\matK}_{ti}\mid\rho)=\tfrac{1}{m}(1+4\rho^2+3\rho^4)$.
Using the law of total variance gives
\(
\sigma^2=\E[\Var(\hat{\matK}_{ti}\mid\rho)]+\Var(\rho^2)
=\Theta(\tfrac{1}{m})+\Theta(\tfrac{1}{d^2})
\),
where $\Var(\rho^2)=\Theta(1/d^2)$ follows from $\E[\rho^2]=1/d$ and $\E[\rho^4]=3/(d(d+2))$.
\end{proof}

\begin{lemma}[Gaussian-chaos entry concentration, centered at $\rho^2$]\label{lem:chaos-centered-conditional}
Let $\mathbf{x},\mathbf{y} \in \R^d$ be deterministic unit vectors and set $\rho \coloneqq \langle \mathbf{x}, \mathbf{y}\rangle$.
Let $\{\va_{r},\vb_{r}\}_{r=1}^m$ be i.i.d.\ with $\va_{r},\vb_{r} \sim \cN(0,I_{d})$, and define the bilinear random-feature kernel
\[
    \hat K(\mathbf{x},\mathbf{y}) \coloneqq \frac1m \sum_{r=1}^m (\va_{r}^\top \mathbf{x})(\va_{r}^\top \mathbf{y})(\vb_{r}^\top \mathbf{x})(\vb_{r}^\top \mathbf{y}).
\]
Then there exists a universal constant $C>0$ such that, for all $u \ge 0$,
\[
    \Pbb\Big(
        \big|\hat K(\mathbf{x},\mathbf{y}) - \rho^2\big|
        \ge
        C\Big(\sqrt{\tfrac{u}{m}} + \tfrac{u^2}{m}\Big)
    \Big)
    \le 2e^{-u}.
\]
\end{lemma}
\begin{proof}
Fix $\mathbf{x},\mathbf{y}$ and write
\[
    Z_{r} \coloneqq (\va_{r}^\top \mathbf{x})(\va_{r}^\top \mathbf{y})(\vb_{r}^\top \mathbf{x})(\vb_{r}^\top \mathbf{y}) - \rho^2,
    \qquad r\in[m].
\]
Then $\{Z_{r}\}_{r=1}^m$ are i.i.d.\ and mean-zero.
Moreover, $Z_{r}$ is a centered degree-$4$ polynomial in jointly Gaussian random variables with bounded covariance, so $\|Z_{r}\|_{L_2}\le C$; by Gaussian hypercontractivity (Theorem~\ref{thm:gaussian-hypercontractivity} with $k=4$), $\|Z_{r}\|_{L_p}\le (p-1)^2\|Z_{r}\|_{L_2}\le C'p^2$ for all $p\ge2$. A standard moment-to-tail conversion then yields a sub-Weibull$(\tfrac12)$ tail: there exists a universal constant $C_{0}>0$ such that, for all $u\ge 0$,
\[
    \Pbb\Big( |Z_{r}| \ge C_{0}(\sqrt{u} + u^2)\Big) \le 2e^{-u}.
\]
Applying Lemma~\ref{lem:weighted-bern-subweib} with weights $w_{r} \equiv 1/m$ (so that $\|\mathbf{w}\|_{2} = m^{-1/2}$ and $\|\mathbf{w}\|_\infty = m^{-1}$) yields the claim.
\end{proof}

\begin{lemma}[Gaussian-chaos entry concentration, centered at $\mu$]\label{lem:chaos-centered-mu}
Let $\mathbf{x},\mathbf{y} \stackrel{\mathrm{i.i.d.}}{\sim} \mathrm{Unif}(\Sph)$ be independent of $\{\va_{r},\vb_{r}\}_{r=1}^m$, and set $\mu \coloneqq 1/d$ and $\rho\coloneqq \langle \mathbf{x},\mathbf{y}\rangle$.
Fix $\delta\in(0,1)$ and define $L \coloneqq \log(6/\delta)$.
Assume
\begin{equation}\label{eq:m-assumption}
    m \le \frac{d^2}{L}.
\end{equation}
Then there exists a universal constant $C>0$ such that
\[
    \Pbb\Big(
        \big|\hat K(\mathbf{x},\mathbf{y}) - \mu\big|
        \le
        C\Big(
            \sqrt{\tfrac{L}{m}} + \tfrac{L^2}{m} + \sqrt{\tfrac{L}{d^2}}
        \Big)
    \Big)
    \ge 1-\delta.
\]
\end{lemma}
\begin{proof}
By the triangle inequality,
\[
    |\hat K(\mathbf{x},\mathbf{y}) - \mu|
    \le
    |\hat K(\mathbf{x},\mathbf{y}) - \rho^2| + |\rho^2 - \mu|.
\]

\paragraph{Control random-feature fluctuation.}
Condition on $\mathbf{x},\mathbf{y}$ and apply Lemma~\ref{lem:chaos-centered-conditional} with $u=L$ to obtain
\[
    \Pbb\Big(
        |\hat K(\mathbf{x},\mathbf{y}) - \rho^2|
        \le
        C\Big(\sqrt{\tfrac{L}{m}} + \tfrac{L^2}{m}\Big)
        \,\Big|\, \mathbf{x},\mathbf{y}
    \Big)
    \ge 1-\delta/3
\]
for a suitable universal constant $C>0$.

\paragraph{Control geometric fluctuation.}
Recall that if $\mathbf{x},\mathbf{y} \stackrel{\text{i.i.d.}}{\sim} \mathrm{Unif}(\Sph)$, then $\rho \stackrel{d}{=} g_{1}/\|g\|_{2}$ for $g \sim \cN(0,I_{d})$.
Let $X \coloneqq \|g\|_{2}^2 \sim \chi^2_{d}$, so that $\rho^2 = g_1^2/X$.
Since $g_1^2 \sim \chi^2_1$ and $X \sim \chi^2_d$, Theorem~\ref{thm:laurent-massart} implies that, each with probability at least $1-\delta/3$,
\[
    |g_1^2 - 1| \le C(\sqrt{L} + L),
    \qquad
    |X - d| \le C(\sqrt{dL} + L)
\]
for a universal constant $C>0$.
In particular, whenever $d \ge C'L$ for a suitable absolute constant $C'$, the second estimate forces $X \ge d/2$; when $d \lesssim L$ the final bound of this step is vacuous up to constants and can be absorbed by enlarging $C$. So assume $X \ge d/2$. Writing $\rho^2 - \mu = \dfrac{g_1^2 d - X}{Xd}$ and using $X \ge d/2$ together with the triangle inequality $|g_1^2 d - X| \le d\,|g_1^2 - 1| + |X - d|$,
\[
    |\rho^2 - \mu|
    =
    \frac{|g_1^2 d - X|}{Xd}
    \le
    \frac{2}{d}\,|g_1^2 - 1| + \frac{2}{d^2}\,|X - d|.
\]
Substituting the two chi-square estimates and using $d \ge 1$,
\[
    |\rho^2 - \mu|
    \le
    \frac{C}{d}(\sqrt{L} + L) + \frac{C}{d^2}(\sqrt{dL} + L)
    \le
    \frac{C}{d}\big(\sqrt{L} + L\big).
\]
Thus, on the intersection of the two chi-square events (of probability at least $1-2\delta/3$),
\[
    |\rho^2 - \mu|
    \le
    \frac{C}{d}\big(\sqrt{L} + L\big).
\]

\paragraph{Combine and simplify.}
Taking a union bound over both gives, with probability at least $1-\delta$,
\[
    |\hat K(\mathbf{x},\mathbf{y}) - \mu|
    \le
    C\Big(
        \sqrt{\tfrac{L}{m}} + \tfrac{L^2}{m} + \frac{\sqrt{L}}{d} + \frac{L}{d}
    \Big).
\]
Finally, under \eqref{eq:m-assumption}, we have $L/d \le \sqrt{L/m}$, so the $L/d$ term is absorbed, yielding the stated bound.
\end{proof}


\begin{lemma}[Centered column-energy concentration \(E_K^\circ\) (isotropic keys, bilinear kernel; sharpened)]
\label{lem:EK-iso-keys}
Assume keys \(\vk_1,\dots,\vk_F\) are i.i.d. uniform on \(\Sph\), independent of
the bilinear random features from \cref{lem:bilinear-sketched-k2}. Let
\[
\mu\defeq \frac1d,
\qquad
\sigma^2\defeq
\E[(\hat{\matK}_{ti}-\mu)^2]
=
\Theta\left(\frac1{d^2}+\frac1m\right),
\qquad t\neq i.
\]
Fix \(\delta\in(0,1)\), and define
\[
L\defeq \log\left(\frac{C_0F^2}{\delta}\right).
\]
Assume
\[
m\le \frac{d^2}{L},
\qquad
L^3\le c_0\sigma^2m^2.
\]
Then, for \(c_0>0\) sufficiently small and \(C_0\) sufficiently large, with
probability at least \(1-\delta\),
\[
E_K^\circ
\le
C\sigma^2\bigl((F-2)+L\bigr).
\]
Equivalently,
\[
\sqrt{E_K^\circ}
\le
C\sigma\sqrt{(F-2)+L}.
\]
In particular, if \(F-2\ge L\), then
\[
\sqrt{E_K^\circ}
\le
C\sigma\sqrt{F-2}.
\]
\end{lemma}

\begin{proof}
Assume
\(F\ge3\). For each \(i\in[F]\), define
\[
\mathcal C_i
\defeq
\sum_{t\neq i}(\hat{\matK}_{ti}-\mu)^2.
\]
For each \(i\neq j\),
\[
E_K^\circ(i,j)
=
\sum_{t\notin\{i,j\}}(\hat{\matK}_{ti}-\mu)^2
\le
\mathcal C_i.
\]
Thus
\[
E_K^\circ\le \max_i\mathcal C_i.
\]

Fix \(i\), and set \(q_i=\vk_i\). For \(q\in\Sph\), define \(A(q)\) as in
Lemma~\ref{lem:bilinear-column-matrix-regularity}. Then, for every \(x\in\Sph\),
\[
x^\top A(q)x
=
\frac1m\sum_{r=1}^m
(\va_r^\top x)(\va_r^\top q)(\vb_r^\top x)(\vb_r^\top q)
=
\hat K(x,q).
\]
Therefore
\[
\hat{\matK}_{ti}=\vk_t^\top A(q_i)\vk_t,
\qquad t\neq i.
\]
Define
\[
\tau_i\defeq \frac{\operatorname{tr}A(q_i)}{d},
\qquad
B_i\defeq A(q_i)-\tau_iI_d,
\qquad
\beta_i\defeq \tau_i-\mu.
\]
Then \(\operatorname{tr}B_i=0\), and
\[
\hat{\matK}_{ti}-\mu
=
\vk_t^\top B_i\vk_t+\beta_i.
\]
Let
\[
\mathcal G_i
\defeq
\sigma\left(\vk_i,\{(\va_r,\vb_r)\}_{r=1}^m\right).
\]
Conditional on \(\mathcal G_i\), the vectors \(\{\vk_t:t\neq i\}\) are independent
uniform spherical vectors. Define
\[
v_i\defeq \frac{2\|B_i\|_F^2}{d(d+2)}+\beta_i^2,
\qquad
r_i\defeq \frac{\|B_i\|_{\mathrm{op}}}{d}+|\beta_i|.
\]
By Lemma~\ref{lem:bilinear-column-parameters}, with probability at least
\(1-\delta/4\), simultaneously for every \(i\in[F]\),
\[
v_i\le C\sigma^2,
\qquad
r_i^2L^2\le C\sigma^2L.
\]
On this event, Lemma~\ref{lem:spherical-quad-square-sum}, applied conditionally
with \(n=F-1\) and \(s=L\), gives
\[
\Pbb\left(
\mathcal C_i
>
C\left((F-1)\sigma^2+\sigma^2L\right)
\,\middle|\,
\mathcal G_i
\right)
\le e^{-L},
\]
because \(L\ge \log(e(F-1))\) after increasing \(C_0\). Union bounding over
\(i\in[F]\),
\[
Fe^{-L}
=
F\frac{\delta}{C_0F^2}
\le
\delta/4
\]
for \(C_0\) sufficiently large. Combining the conditional square-sum event with
the parameter event yields, with probability at least \(1-\delta\),
\[
\max_i\mathcal C_i
\le
C\sigma^2\bigl((F-1)+L\bigr).
\]
Since \(F\ge3\), \(F-1\le2(F-2)\), so
\[
E_K^\circ
\le
C\sigma^2\bigl((F-2)+L\bigr).
\]
Taking square roots proves
\[
\sqrt{E_K^\circ}
\le
C\sigma\sqrt{(F-2)+L}.
\]
If \(F-2\ge L\), then \((F-2)+L\le2(F-2)\), giving
\[
\sqrt{E_K^\circ}
\le
C\sigma\sqrt{F-2}.
\]
\end{proof}



\begin{lemma}[Concentration of the effective coupling $\kappa_{\mathrm{eff}}^\circ$ (isotropic keys, bilinear kernel)]
\label{lem:keff-iso-keys}
Assume the setting of Theorem~\ref{app:regime2}: isotropic keys, bilinear kernel, and deterministic values/codes.
Let $\mu=1/d$ and $\sigma^2=\Theta(\tfrac{1}{d^2}+\tfrac{1}{m})$.
Fix $\delta\in(0,1)$ and let $L\defeq \log(\tfrac{C_{0} F^2}{\delta})$.
Assume additionally that $m\le d^2/L$.
Define
\[
\kappa_{\mathrm{eff}}^\circ
\ \defeq\
\max_{i\neq j}\frac{|\ip{\vX^{\circ(ij)}}{\vY^{(ij)}}|}{\|\vY^{(ij)}\|_2}
\]
and
\[
L_{v}\ \defeq\ \max_{i\neq j}\frac{\|\vY^{(ij)}\|_{1}^2}{\|\vY^{(ij)}\|_{2}^2}
\]
Then with probability at least $1-\delta$ (over keys and features),
\[
\kappa_{\mathrm{eff}}^\circ
\ \le\
C\left(\sigma\sqrt{L}+\frac{L^2}{m}\right)\sqrt{L_v}.
\]
\end{lemma}

\begin{proof}
Fix a pair $(i,j)$ with $j\neq i$.
Write $\theta_{t}\defeq \hat{\matK}_{ti}-\mu$ for $t\notin\{i,j\}$, so that
$\vX^{\circ(ij)}=(\theta_{t})_{t\notin\{i,j\}}$.
Then
\[
\frac{|\ip{\vX^{\circ(ij)}}{\vY^{(ij)}}|}{\|\vY^{(ij)}\|_2}
\ =\
\frac{\big|\sum_{t\notin\{i,j\}} \theta_{t}\,\vY^{(ij)}_{t}\big|}{\|\vY^{(ij)}\|_2}
\ \le\
\Big(\max_{t\notin\{i,j\}}|\theta_{t}|\Big)\,\frac{\|\vY^{(ij)}\|_1}{\|\vY^{(ij)}\|_2}.
\]
Taking a maximum over $(i,j)$ and using the definition of $L_{v}$ gives
\[
\kappa_{\mathrm{eff}}^\circ
\ \le\
\Big(\max_{p\neq q}|\hat{\matK}_{pq}-\mu|\Big)\,\sqrt{L_v}.
\]
Now apply Lemma~\ref{lem:Koff-iso-keys} with failure probability $\delta$ to bound
\(
\max_{p\neq q}|\hat{\matK}_{pq}-\mu|\le C(\sigma\sqrt{L}+\tfrac{L^2}{m}).
\)
Multiplying by $\sqrt{L_v}$ completes the proof.
\end{proof}



\subsubsection{Additional Concentration Tools}
\label{app:aux-additional-concentration-coupling}

\begin{lemma}[Weighted Bernstein for sub-Weibull$(1/2)$-type tails]
\label{lem:weighted-bern-subweib}
Let $Z_1,\dots,Z_N$ be independent, mean-zero random variables such that
\[
\Pbb\!\left(|Z_r|\ge C_0(\sigma\sqrt{u}+u^2/m)\right)\le 2e^{-u}
\qquad\text{for all }u\ge0,\ r\in[N].
\]
Then there exist absolute constants $C,c>0$ such that for every deterministic
$w\in\R^N$ and every $t\ge0$,
\[
\Pbb\!\left(
\Big|\sum_{r=1}^N w_r Z_r\Big|
\ge
C\Big(\sigma\|w\|_2\sqrt{t}+\frac{t^2}{m}\|w\|_\infty\Big)
\right)
\le 2e^{-ct}.
\]
Consequently, for $u=\log(C_1/\delta)$, with probability at least $1-\delta$,
\[
\Big|\sum_{r=1}^N w_r Z_r\Big|
\le
C_2\Big(\sigma\|w\|_2\sqrt{u}+\frac{u^2}{m}\|w\|_\infty\Big).
\]
\end{lemma}
\begin{proof}
The assumed tail bound is equivalent up to absolute constants to the statement
that each $Z_r$ is sub-Weibull of order $\alpha=\tfrac12$ with generalized
Bernstein--Orlicz parameters $\nu\asymp \sigma$ and $L\asymp (\sigma m)^{-1}$.
Applying Theorem~2.4 of \citep{BongKuchibhotla2023} to the weighted sum
$\sum_{r=1}^N w_r Z_r$ gives
\[
\Pbb\!\left(
\Big|\sum_{r=1}^N w_r Z_r\Big|
\ge
C\Big(\sigma\|w\|_2\sqrt{t}+\sigma L\, t^2 \|w\|_\infty\Big)
\right)
\le 2e^{-ct}.
\]
Since $\sigma L \asymp 1/m$, this becomes
\[
\Pbb\!\left(
\Big|\sum_{r=1}^N w_r Z_r\Big|
\ge
C\Big(\sigma\|w\|_2\sqrt{t}+\frac{t^2}{m}\|w\|_\infty\Big)
\right)
\le 2e^{-ct}.
\]
Setting $t=\log(C_1/\delta)$ yields the stated high-probability bound.
\end{proof}



\begin{theorem}[Gaussian hypercontractivity; Nelson--Gross]
\label{thm:gaussian-hypercontractivity}
Let \(G\sim\mathcal N(0,I_n)\), and let \(P(G)\) be a polynomial of degree at
most \(k\) in the standard Gaussian variables. Then, for every \(q\ge2\),
\[
\|P(G)\|_{L_q}
\le
(q-1)^{k/2}\|P(G)\|_{L_2}.
\]
\end{theorem}

\noindent
This is the standard degree-\(k\) polynomial-chaos corollary of the Gaussian
hypercontractivity theorem; see O'Donnell~\cite[Theorem~11.23]{odonnell2014analysis}.

\begin{theorem}[Sub-Weibull average bound from moment growth]
\label{thm:subweibull-bernstein-average}
Let \(W_1,\dots,W_m\) be independent mean-zero random variables. Suppose that,
for some \(\rho\ge1\) and \(K>0\),
\[
\|W_r\|_{L_p}\le Kp^\rho
\qquad\text{for every }p\ge2\text{ and every }r\in[m].
\]
Then there is a constant \(C_\rho>0\), depending only on \(\rho\), such that for
every \(u\ge1\),
\[
\Pbb\left(
\left|
\frac1m\sum_{r=1}^m W_r
\right|
>
C_\rho K\left(
\sqrt{\frac{u}{m}}+
\frac{u^\rho}{m}
\right)
\right)
\le 2e^{-u}.
\]
In particular, if \(W_1,\dots,W_m\) are independent centered degree-\(k\)
Gaussian chaoses with \(k\ge2\) and
\[
\|W_r\|_{L_2}\le K_0,
\]
then
\[
\Pbb\left(
\left|
\frac1m\sum_{r=1}^m W_r
\right|
>
C_kK_0\left(
\sqrt{\frac{u}{m}}+
\frac{u^{k/2}}{m}
\right)
\right)
\le 2e^{-u}.
\]
\end{theorem}

\begin{proof}
Set \(\alpha\defeq1/\rho\in(0,1]\). By a standard argument, the moment growth
condition implies the Orlicz-norm bound
\[
\|W_r\|_{\psi_\alpha}
\defeq
\inf\left\{c>0:\E\exp\left((|W_r|/c)^\alpha\right)\le2\right\}
\le
C_\rho K
\qquad\text{for every }r\in[m]:
\]
indeed, monotonicity of \(p\mapsto\|W_r\|_{L_p}\) and the assumption give
\(\|W_r\|_{L_{\alpha j}}\le C_\rho Kj^{1/\alpha}\) for every integer \(j\ge1\),
and expanding the exponential and using \(j!\ge(j/e)^j\) shows
\(\E\exp\left((|W_r|/(A_\rho K))^\alpha\right)\le2\) for \(A_\rho\) sufficiently
large.

Now apply the generalized Bernstein--Orlicz inequality ~\cite[Theorem~3.1]{kuchibhotla2022moving} to
\(Z\defeq\frac1m\sum_{r=1}^mW_r\), with weights \(a_r=1/m\) and
\(b_r\defeq a_r\|W_r\|_{\psi_\alpha}\), so that
\(\|b\|_2\le C_\rho K/\sqrt m\) and \(\|b\|_\infty\le C_\rho K/m\).
The defining tail property of their generalized Bernstein--Orlicz norm yields,
for every \(u\ge1\),
\[
\Pbb\left(
|Z|
>
C_\alpha\left(
\|b\|_2\sqrt u
+
\|b\|_\infty u^{1/\alpha}
\right)
\right)
\le
2e^{-u}.
\]
Substituting the bounds on \(\|b\|_2\) and \(\|b\|_\infty\) and using
\(1/\alpha=\rho\) gives the stated tail bound.

Finally, if \(W_r\) is a centered degree-\(k\) Gaussian chaos with
\(\|W_r\|_{L_2}\le K_0\), Theorem~\ref{thm:gaussian-hypercontractivity} gives
\[
\|W_r\|_{L_p}
\le
(p-1)^{k/2}\|W_r\|_{L_2}
\le
C_kK_0p^{k/2},
\qquad p\ge2.
\]
Apply the first part with \(\rho=k/2\) and \(K=C_kK_0\).
\end{proof}
\begin{theorem}[Operator norm from bilinear forms on nets]
\label{thm:operator-norm-net}
Let \(A\in\mathbb R^{n\times m}\), let \(0<\varepsilon<1/2\), and let
\(\mathcal N\subset \mathbb S^{n-1}\), \(\mathcal M\subset\mathbb S^{m-1}\)
be finite \(\varepsilon\)-nets. Then
\[
\max_{x\in\mathcal N,\;y\in\mathcal M}|x^\top Ay|
\le
\|A\|_{\mathrm{op}}
\le
\frac{1}{1-2\varepsilon}
\max_{x\in\mathcal N,\;y\in\mathcal M}|x^\top Ay|.
\]
In particular, for \(\varepsilon=1/4\),
\[
\|A\|_{\mathrm{op}}
\le
2\max_{x\in\mathcal N,\;y\in\mathcal M}|x^\top Ay|.
\]
Moreover, for every \(d\ge 1\), the sphere \(\mathbb S^{d-1}\) admits a
\(1/4\)-net of cardinality at most \(9^d\). Hence one may choose such nets with
\[
|\mathcal N|\le 9^n
\qquad\text{and}\qquad
|\mathcal M|\le 9^m .
\]
\end{theorem}

\noindent
This is Vershynin~\cite[Lemma~4.4.2 and Corollary~4.2.11]{vershynin2026high}.

\begin{theorem}[Laurent--Massart weighted chi-square tail]
\label{thm:laurent-massart}
Let \(Y_1,\dots,Y_D\) be independent standard Gaussian random variables, and let
\(a_1,\dots,a_D\ge0\). Define
\[
Z\defeq \sum_{i=1}^D a_i(Y_i^2-1).
\]
Then, for every \(x>0\),
\[
\Pbb\left(
Z\ge 2\|a\|_2\sqrt x+2\|a\|_\infty x
\right)
\le e^{-x},
\]
and
\[
\Pbb\left(
Z\le -2\|a\|_2\sqrt x
\right)
\le e^{-x}.
\]
In particular, if \(Q\sim\chi^2_D\), then
\[
\Pbb\left(Q-D\ge 2\sqrt{Dx}+2x\right)
\le e^{-x},
\]
and
\[
\Pbb\left(D-Q\ge 2\sqrt{Dx}\right)
\le e^{-x}.
\]
\end{theorem}

\noindent
This is Laurent--Massart~\cite[Lemma~1]{laurent2000adaptive}.




\begin{lemma}[Weighted product-Gaussian chaos]
\label{lem:weighted-product-gaussian}
Let \((g_r,h_r)_{r=1}^m\) be independent pairs of independent standard
Gaussians, and let \(a=(a_1,\dots,a_m)\in\R^m\) be deterministic. Then, for every
\(t\ge0\),
\[
\Pbb\left(
\left|
\sum_{r=1}^m a_rg_rh_r
\right|
>
C\left(\|a\|_2\sqrt t+\|a\|_\infty t\right)
\right)
\le 2e^{-t}.
\]
\end{lemma}

\begin{proof}
For independent \(g,h\sim N(0,1)\), conditioning on \(g\) gives
\[
\E e^{\lambda gh}
=
\E e^{\lambda^2g^2/2}
=
(1-\lambda^2)^{-1/2},
\qquad |\lambda|<1.
\]
Therefore, for \(|\lambda|\le c/\|a\|_\infty\),
\[
\log\E\exp\left(\lambda\sum_{r=1}^m a_rg_rh_r\right)
=
-\frac12\sum_{r=1}^m\log(1-\lambda^2a_r^2)
\le
C\lambda^2\|a\|_2^2.
\]
If \(a=0\), the claim is trivial. Otherwise, Chernoff's bound gives
\[
\Pbb\left(
\sum_{r=1}^m a_rg_rh_r
>
\nu
\right)
\le
\exp\left(-\lambda\nu+C\lambda^2\|a\|_2^2\right)
\]
for every \(0\le\lambda\le c/\|a\|_\infty\). Optimizing with
\[
\lambda
=
 c\min\left\{\frac{\nu}{\|a\|_2^2},\frac1{\|a\|_\infty}\right\}
\]
gives
\[
\Pbb\left(
\sum_{r=1}^m a_rg_rh_r
>
C(\|a\|_2\sqrt t+\|a\|_\infty t)
\right)
\le e^{-t}.
\]
Apply the same argument to \(-\sum_r a_rg_rh_r\) and union bound.
\end{proof}

\begin{lemma}[Square-sum concentration from a two-level tail]
\label{lem:square-sum-two-level}
Let \(X_1,\dots,X_n\) be independent random variables. Assume that, for constants
\(v>0\), \(r>0\), and \(C_0\ge1\),
\[
\E X_i^2\le v
\]
and, for every \(u\ge1\),
\[
\Pbb\left(
|X_i|>C_0(\sqrt{vu}+ru)
\right)
\le 2e^{-u}.
\]
Assume also that
\[
r^2\le C_0v.
\]
Then, for every \(s\ge1\),
\[
\Pbb\left(
\sum_{i=1}^nX_i^2
>
C\left(nv+vs+r^2s^2\right)
\right)
\le Ce^{-s},
\]
where \(C\) depends only on \(C_0\).
\end{lemma}

\begin{proof}
We first reduce the problem to sums of variables with exponential tails. For each
\(i\), let \(R_i\) have the same distribution as \(|X_i|\), with the variables
\(R_1,\dots,R_n\) independent. Define
\[
T_i
\defeq
\inf\left\{u\ge1:
R_i\le C_0(\sqrt{vu}+ru)
\right\}.
\]
Then, by definition,
\[
R_i\le C_0(\sqrt{vT_i}+rT_i).
\]
Moreover, for every \(u\ge1\), if \(T_i>u\), then
\[
R_i>C_0(\sqrt{vu}+ru),
\]
and hence the assumed tail bound gives
\[
\Pbb(T_i>u)\le 2e^{-u}.
\]
Since the desired estimate depends only on the marginal laws and independence, we may
work with this coupled representation and write, after changing the absolute constant,
\[
|X_i|
\le
C(\sqrt{vT_i}+rT_i),
\qquad
\Pbb(T_i>u)\le2e^{-u},\quad u\ge1.
\]

Squaring the domination gives
\[
X_i^2
\le
C\left(vT_i+r^2T_i^2\right).
\]
Hence
\[
\sum_{i=1}^n X_i^2
\le
C\left(
 v\sum_{i=1}^nT_i+r^2\sum_{i=1}^nT_i^2
\right).
\]
It remains to control the two sums involving \(T_i\).

The tail bound \(\Pbb(T_i>u)\le2e^{-u}\) implies, for every \(p\ge2\),
\[
\|T_i\|_{L_p}
\le Cp,
\qquad
\|T_i^2\|_{L_p}
=
\|T_i\|_{L_{2p}}^2
\le Cp^2.
\]
Consequently,
\[
\|T_i-\E T_i\|_{L_p}\le Cp,
\qquad
\|T_i^2-\E T_i^2\|_{L_p}\le Cp^2.
\]
Also \(\E T_i\le C\) and \(\E T_i^2\le C\). Applying
Theorem~\ref{thm:subweibull-bernstein-average} to the centered variables
\(T_i-\E T_i\) with \(\rho=1\) gives
\[
\Pbb\left(
\sum_{i=1}^nT_i>C(n+s)
\right)
\le Ce^{-s}.
\]
Similarly, applying Theorem~\ref{thm:subweibull-bernstein-average} to
\(T_i^2-\E T_i^2\) with \(\rho=2\) gives
\[
\Pbb\left(
\sum_{i=1}^nT_i^2>C(n+s^2)
\right)
\le Ce^{-s}.
\]
On the intersection of these two events,
\[
\sum_{i=1}^n X_i^2
\le
C\left(v(n+s)+r^2(n+s^2)\right).
\]
Since \(r^2\le C_0v\), the term \(r^2n\) is absorbed by \(vn\). Therefore
\[
\sum_{i=1}^n X_i^2
\le
C\left(nv+vs+r^2s^2\right)
\]
with probability at least \(1-Ce^{-s}\).
\end{proof}

\begin{lemma}[Spherical quadratic-form tail]
\label{lem:spherical-quadratic-tail}
Let \(x\sim\mathrm{Unif}(\Sph)\), and let \(B\in\R^{d\times d}\) be
symmetric and trace-free. Then, for every \(u\ge1\),
\[
\Pbb\left(
|x^\top Bx|
>
C\left(
\frac{\|B\|_F}{d}\sqrt u+
\frac{\|B\|_{\mathrm{op}}}{d}u
\right)
\right)
\le
2e^{-u}.
\]
\end{lemma}

\begin{proof}
Write
\[
x=\frac{g}{\|g\|_2},
\qquad
 g\sim\mathcal N(0,I_d).
\]
Diagonalize \(B=Q\operatorname{diag}(\lambda_1,\dots,\lambda_d)Q^\top\). By
rotational invariance,
\[
g^\top Bg
\stackrel{d}{=}
\sum_{i=1}^d\lambda_i g_i^2.
\]
Since \(\operatorname{tr}B=\sum_i\lambda_i=0\),
\[
g^\top Bg
=
\sum_{i=1}^d\lambda_i(g_i^2-1).
\]
Apply Theorem~\ref{thm:laurent-massart} to the positive and negative parts of the
weights \(\lambda_i\). Equivalently, apply it to the two weighted sums associated
with \(\lambda_i^+\) and \(\lambda_i^-\). This gives, for all \(u\ge1\),
\[
\Pbb\left(
|g^\top Bg|
>
C\left(\|B\|_F\sqrt u+\|B\|_{\mathrm{op}}u\right)
\right)
\le
2e^{-u}.
\]
Choose a numerical constant \(a\ge1\), to be fixed below. Applying the above
Gaussian quadratic-form bound with \(au\) in place of \(u\), we get
\[
\Pbb\left(
|g^\top Bg|
>
C_a\left(\|B\|_F\sqrt u+\|B\|_{\mathrm{op}}u\right)
\right)
\le 2e^{-au}.
\]
Also, by the lower-tail part of Theorem~\ref{thm:laurent-massart} applied to
\(\|g\|_2^2\sim\chi_d^2\),
\[
\Pbb\left(\|g\|_2^2<d/2\right)
\le e^{-cd}.
\]
If \(u\le c_1d\), with \(c_1>0\) chosen small enough, then
\(e^{-cd}\le e^{-2u}\). On the event \(\|g\|_2^2\ge d/2\),
\[
|x^\top Bx|
=
\frac{|g^\top Bg|}{\|g\|_2^2}
\le
C_a\left(
\frac{\|B\|_F}{d}\sqrt u+
\frac{\|B\|_{\mathrm{op}}}{d}u
\right).
\]
Thus, for \(u\le c_1d\), the failure probability is at most
\(2e^{-au}+e^{-2u}\). Choosing \(a\) large enough and enlarging the constant in
the threshold gives failure probability at most \(2e^{-u}\).
If \(u>c_1d\), then the same inequality holds deterministically after
enlarging \(C\), because
\[
|x^\top Bx|\le\|B\|_{\mathrm{op}}
\le
C\frac{\|B\|_{\mathrm{op}}}{d}u.
\]
This proves the claim.
\end{proof}

\begin{lemma}[Spherical quadratic square-sum bound]
\label{lem:spherical-quad-square-sum}
Let \(x_1,\dots,x_n\) be i.i.d. uniform on \(\Sph\subset\R^d\). Let
\(B\in\R^{d\times d}\) be deterministic, symmetric, and trace-free, and let
\(\beta\in\R\). Define
\[
X_\ell\defeq x_\ell^\top Bx_\ell+\beta.
\]
Set
\[
v\defeq \frac{2\|B\|_F^2}{d(d+2)}+\beta^2,
\qquad
r\defeq \frac{\|B\|_{\mathrm{op}}}{d}+|\beta|.
\]
Then, for every \(s\ge1\),
\[
\Pbb\left(
\sum_{\ell=1}^nX_\ell^2
>
C\left(nv+vs+r^2s^2\right)
\right)
\le
e^{-s}.
\]
\end{lemma}

\begin{proof}
For \(x\sim\mathrm{Unif}(\Sph)\), isotropy gives
\[
\E[x^\top Bx]=\frac{\operatorname{tr}B}{d}=0.
\]
The fourth-moment identity for the sphere gives
\[
\E[(x^\top Bx)^2]
=
\frac{(\operatorname{tr}B)^2+2\|B\|_F^2}{d(d+2)}
=
\frac{2\|B\|_F^2}{d(d+2)}.
\]
Therefore
\[
\E X_\ell^2=v.
\]
Moreover,
\[
\left(\frac{\|B\|_{\mathrm{op}}}{d}\right)^2
\le
\frac{\|B\|_F^2}{d^2}
\le
C\frac{2\|B\|_F^2}{d(d+2)}
\le
Cv,
\]
and \(\beta^2\le v\), so
\[
r^2\le Cv.
\]
By Lemma~\ref{lem:spherical-quadratic-tail}, for every \(u\ge1\),
\[
\Pbb\left(
|x^\top Bx|
>
C\left(
\frac{\|B\|_F}{d}\sqrt u+
\frac{\|B\|_{\mathrm{op}}}{d}u
\right)
\right)
\le 2e^{-u}.
\]
Since \(\|B\|_F/d\le C\sqrt v\), \(\|B\|_{\mathrm{op}}/d\le r\), and
\(|\beta|\le\sqrt v\), this implies
\[
\Pbb\left(
|X_\ell|>C(\sqrt{vu}+ru)
\right)
\le 2e^{-u}.
\]
Lemma~\ref{lem:square-sum-two-level} applies because \(r^2\le Cv\). It gives
\[
\Pbb\left(
\sum_{\ell=1}^nX_\ell^2
>
C\left(nv+vs+r^2s^2\right)
\right)
\le Ce^{-s}.
\]
Replacing \(s\) in the equation above by \(s+c_0\), for a sufficiently large
absolute constant \(c_0\), and enlarging \(C\), changes the right-hand side to
\(e^{-s}\) and leaves the threshold in the same form. This proves the
claim.
\end{proof}

\begin{lemma}[Scalar weight regularity for one bilinear column]
\label{lem:scalar-weight-regularity}
Let \(\alpha_1,\dots,\alpha_m,\beta_1,\dots,\beta_m\) be independent standard
Gaussians, and define
\[
\gamma_r=\alpha_r\beta_r,
\qquad
\eta_r=\alpha_r\beta_r^2,
\qquad
\zeta_r=\alpha_r^2\beta_r.
\]
If \(m\ge C_*L^3\), then, with probability at least \(1-Ce^{-L}\),
\[
\left|
\frac1m\sum_{r=1}^m(\alpha_r^2\beta_r^2-1)
\right|
\le
C\left(\sqrt{\frac Lm}+\frac{L^2}{m}\right),
\]
\[
\sum_{r=1}^m\gamma_r^2\le Cm,
\qquad
\sum_{r=1}^m\eta_r^2\le Cm,
\qquad
\sum_{r=1}^m\zeta_r^2\le Cm,
\]
and
\[
\|\gamma\|_\infty^2L\le Cm.
\]
\end{lemma}

\begin{proof}
The variable \(\alpha^2\beta^2-1\) is a centered degree-four Gaussian polynomial
with bounded \(L_2\)-norm. By Theorem~\ref{thm:gaussian-hypercontractivity},
\[
\|\alpha^2\beta^2-1\|_{L_p}
\le
Cp^2,
\qquad p\ge2.
\]
Applying Theorem~\ref{thm:subweibull-bernstein-average} with \(\rho=2\) gives
\[
\left|
\frac1m\sum_{r=1}^m(\alpha_r^2\beta_r^2-1)
\right|
\le
C\left(\sqrt{\frac Lm}+\frac{L^2}{m}\right)
\]
with probability at least \(1-2e^{-L}\).

Next,
\[
\E\gamma_r^2=1,
\qquad
\E\eta_r^2=3,
\qquad
\E\zeta_r^2=3.
\]
The centered variables
\[
\gamma_r^2-1,
\qquad
\eta_r^2-3,
\qquad
\zeta_r^2-3
\]
are centered Gaussian polynomials of degrees \(4,6,6\), respectively, with bounded
\(L_2\)-norms. Applying Theorem~\ref{thm:subweibull-bernstein-average} with
\(\rho=2\) for \(\gamma_r^2-1\) and \(\rho=3\) for the degree-six terms gives
\[
\frac1m\sum_{r=1}^m\gamma_r^2\le C,
\qquad
\frac1m\sum_{r=1}^m\eta_r^2\le C,
\qquad
\frac1m\sum_{r=1}^m\zeta_r^2\le C
\]
with probability at least \(1-Ce^{-L}\), provided \(m\ge C_*L^3\).

It remains to control \(\|\gamma\|_\infty\). Since
\[
2|\alpha\beta|\le \alpha^2+\beta^2,
\]
and \(\alpha^2+\beta^2\sim\chi^2_2\), there is an absolute constant \(c>0\) such
that
\[
\Pbb(|\alpha\beta|>u)\le 2e^{-cu}.
\]
Thus
\[
\Pbb\left(\|\gamma\|_\infty^2L>Cm\right)
\le
2m\exp\left(-c\sqrt{\frac{Cm}{L}}\right).
\]
Since \(m\ge C_*L^3\), write \(m=L^3y\) with \(y\ge C_*\). Then
\[
\sqrt{\frac{m}{L}}=L\sqrt y,
\]
whereas
\[
L+\log(2m)+1
=
L+\log(2L^3y)+1
\le
C(L+\log y)
\le
CL\sqrt y.
\]
Choosing the absolute constants \(C_*\) and \(C\) sufficiently large therefore ensures
\[
c\sqrt{\frac{Cm}{L}}
\ge
L+\log(2m)+1.
\]
Therefore
\[
2m\exp\left(-c\sqrt{\frac{Cm}{L}}\right)
\le e^{-L}.
\]
A union bound over the above events completes the proof.
\end{proof}

\begin{lemma}[Weighted Gaussian product bounds]
\label{lem:weighted-gaussian-product}
Let \(U,V\in\R^{n\times m}\) have independent standard Gaussian entries, and let
\(\Gamma=\operatorname{diag}(\gamma_1,\dots,\gamma_m)\) be deterministic. Then,
for every \(s\ge1\), with probability at least \(1-4e^{-s}\),
\[
\|U\Gamma V^\top\|_{\mathrm{op}}
\le
C\left(
\|\gamma\|_2\sqrt{n+s}+\|\gamma\|_\infty(n+s)
\right),
\]
and
\[
\left|\operatorname{tr}(U\Gamma V^\top)\right|
\le
C\left(\sqrt n\,\|\gamma\|_2\sqrt s+\|\gamma\|_\infty s\right).
\]
Furthermore, if
\[
\|\gamma\|_2^2\le C_0m,
\qquad
\|\gamma\|_\infty^2s\le C_0m,
\qquad
n\ge s,
\]
then, with probability at least \(1-6e^{-s}\),
\[
\|U\Gamma V^\top\|_F^2\le Cn^2m.
\]
\end{lemma}

\begin{proof}
For fixed \(x,y\in\mathbb S^{n-1}\),
\[
x^\top U\Gamma V^\top y
=
\sum_{r=1}^m \gamma_r (u_r^\top x)(v_r^\top y),
\]
where \(u_r,v_r\) are the columns of \(U,V\). The variables
\((u_r^\top x)(v_r^\top y)\) are independent products of independent standard
Gaussians. Lemma~\ref{lem:weighted-product-gaussian} gives, for every \(t\ge0\),
\[
\Pbb\left(
|x^\top U\Gamma V^\top y|
>
C(\|\gamma\|_2\sqrt t+\|\gamma\|_\infty t)
\right)
\le
2e^{-t}.
\]
Take \(1/4\)-nets \(\mathcal N,\mathcal M\) of \(\mathbb S^{n-1}\) with cardinalities
at most \(9^n\). Set \(t=s+2n\log 9\), union bound over
\(\mathcal N\times\mathcal M\), and apply Theorem~\ref{thm:operator-norm-net}.
This proves the operator bound.

For the trace,
\[
\operatorname{tr}(U\Gamma V^\top)
=
\sum_{\ell=1}^n\sum_{r=1}^m\gamma_rU_{\ell r}V_{\ell r}.
\]
This is again a weighted sum of products of independent standard Gaussians, with
weights \(\gamma_r\) repeated \(n\) times. Lemma~\ref{lem:weighted-product-gaussian}
therefore gives the stated trace bound.

For the Frobenius bound, first note
\[
\|U\Gamma\|_F^2
=
\sum_{r=1}^m\gamma_r^2\|u_r\|_2^2.
\]
This is a weighted chi-square variable with mean \(n\|\gamma\|_2^2\le Cnm\). By
Theorem~\ref{thm:laurent-massart}, using
\(\|\gamma\|_\infty^2s\le C_0m\), we get
\[
\|U\Gamma\|_F^2\le Cnm
\]
with probability at least \(1-e^{-s}\).

Condition on \(U,\Gamma\), and set
\[
R\defeq \Gamma U^\top U\Gamma.
\]
Then
\[
\|U\Gamma V^\top\|_F^2
=
\operatorname{tr}(VRV^\top).
\]
Writing the rows of \(V\) as \(g_1,\dots,g_n\in\R^m\),
\[
\operatorname{tr}(VRV^\top)
=
\sum_{\ell=1}^n g_\ell^\top Rg_\ell.
\]
Diagonalize \(R\), and apply Theorem~\ref{thm:laurent-massart}. With conditional
probability at least \(1-e^{-s}\),
\[
\operatorname{tr}(VRV^\top)
\le
n\operatorname{tr}R
+2\sqrt{n\operatorname{tr}(R^2)s}
+2\|R\|_{\mathrm{op}}s.
\]
On \(\{\|U\Gamma\|_F^2\le Cnm\}\),
\[
\operatorname{tr}R=\|U\Gamma\|_F^2\le Cnm,
\]
\[
\|R\|_{\mathrm{op}}\le \operatorname{tr}R\le Cnm,
\]
and
\[
\operatorname{tr}(R^2)
\le
\|R\|_{\mathrm{op}}\operatorname{tr}R
\le
Cn^2m^2.
\]
Since \(n\ge s\), all three terms are bounded by \(Cn^2m\). Thus
\[
\|U\Gamma V^\top\|_F^2\le Cn^2m.
\]
Union bounding the events completes the proof.
\end{proof}

\begin{lemma}[Bilinear one-column matrix regularity]
\label{lem:bilinear-column-matrix-regularity}
Fix \(\delta\in(0,1)\), and define
\[
L\defeq \log\left(\frac{C_0F^2}{\delta}\right).
\]
Assume
\[
m\le \frac{d^2}{L},
\qquad
m\ge C_*L^3.
\]
Let \(\va_1,\vb_1,\dots,\va_m,\vb_m\in\R^d\) be i.i.d.\ standard Gaussian
vectors---the bilinear random features of \cref{lem:bilinear-sketched-k2}.
For \(q\in\Sph\), define
\[
A(q)
\defeq
\frac1{2m}\sum_{r=1}^m
(\va_r^\top q)(\vb_r^\top q)(\va_r\vb_r^\top+\vb_r\va_r^\top),
\]
and
\[
N(q)\defeq A(q)-qq^\top.
\]
Let \(\vk_1,\dots,\vk_F\) be i.i.d. uniform on \(\Sph\), independent of the
features. Then, with probability at least \(1-\delta/4\), simultaneously for all
\(i\in[F]\),
\[
\|N(\vk_i)\|_F^2\le C\frac{d^2}{m},
\]
\[
\|N(\vk_i)\|_{\mathrm{op}}
\le
C\left(\sqrt{\frac dm}+\frac{d}{\sqrt{mL}}\right),
\]
and
\[
|\operatorname{tr}N(\vk_i)|
\le
C\sqrt{\frac{dL}{m}}.
\]
\end{lemma}

\begin{proof}
It suffices to prove the result for a fixed deterministic \(q\in\Sph\) with
failure probability at most \(Ce^{-L}\). Conditional on the keys, the vectors
\(\vk_i\) are deterministic and independent of the feature randomness. A union
bound over \(i\in[F]\) gives failure probability at most
\[
CFe^{-L}=CF\frac{\delta}{C_0F^2}\le \delta/4
\]
for \(C_0\) sufficiently large.

Fix \(q\). By rotational invariance, assume \(q=e_1\). Write
\[
\va_r=(\alpha_r,u_r),
\qquad
\vb_r=(\beta_r,v_r),
\]
where \(\alpha_r,\beta_r\sim N(0,1)\), \(u_r,v_r\sim\mathcal N(0,I_{d-1})\), and
all variables are independent across \(r\). Set \(n=d-1\). The \(r\)-th summand is
\[
S_r
=
\frac12\alpha_r\beta_r(\va_r\vb_r^\top+\vb_r\va_r^\top),
\]
and \(\E S_r=e_1e_1^\top\). In block form,
\[
N(e_1)
=
\begin{pmatrix}
N_{11} & N_{1\perp}^\top\\
N_{1\perp} & N_{\perp\perp}
\end{pmatrix},
\]
where
\[
N_{11}=\frac1m\sum_{r=1}^m(\alpha_r^2\beta_r^2-1),
\]
\[
N_{1\perp}=\frac1{2m}\sum_{r=1}^m(\eta_ru_r+\zeta_rv_r),
\]
and, with \(U,V\in\R^{n\times m}\) having columns \(u_r,v_r\),
\(\Gamma=\operatorname{diag}(\gamma_1,\dots,\gamma_m)\),
\[
N_{\perp\perp}
=
\frac1{2m}(U\Gamma V^\top+V\Gamma U^\top),
\]
with
\[
\gamma_r=\alpha_r\beta_r,
\qquad
\eta_r=\alpha_r\beta_r^2,
\qquad
\zeta_r=\alpha_r^2\beta_r.
\]

Let \(\mathcal G\) be the good event from Lemma~\ref{lem:scalar-weight-regularity}.
On \(\mathcal G\),
\[
|N_{11}|\le C\left(\sqrt{\frac Lm}+\frac{L^2}{m}\right),
\]
\[
\|\gamma\|_2^2\le Cm,
\qquad
\|\gamma\|_\infty^2L\le Cm,
\qquad
\sum_r\eta_r^2\le Cm,
\qquad
\sum_r\zeta_r^2\le Cm.
\]
Also \(\Pbb(\mathcal G^c)\le Ce^{-L}\).

Conditional on the scalar weights,
\[
S_\eta\defeq\sum_r\eta_ru_r
\sim
\mathcal N\left(0,\left(\sum_r\eta_r^2\right)I_n\right),
\qquad
S_\zeta\defeq\sum_r\zeta_rv_r
\sim
\mathcal N\left(0,\left(\sum_r\zeta_r^2\right)I_n\right).
\]
On \(\mathcal G\), both covariance scalars are at most \(Cm\). Therefore
\(\|S_\eta\|_2^2/(Cm)\) and \(\|S_\zeta\|_2^2/(Cm)\) are stochastically dominated,
up to an absolute constant, by \(\chi_n^2\). By the chi-square consequence of
Theorem~\ref{thm:laurent-massart}, with conditional failure probability at most
\(2e^{-L}\),
\[
\|S_\eta\|_2
\le
C\sqrt m(\sqrt n+\sqrt L),
\qquad
\|S_\zeta\|_2
\le
C\sqrt m(\sqrt n+\sqrt L).
\]
Hence
\[
\|N_{1\perp}\|_2
\le
C\sqrt{\frac{n+L}{m}}.
\]
Since \(m\ge C_*L^3\) and \(m\le d^2/L\),
\[
d^2\ge mL\ge C_*L^4,
\]
so \(d\ge cL^2\) and \(n+L\le Cd\). Therefore
\[
\|N_{1\perp}\|_2\le C\sqrt{\frac dm}.
\]

By Lemma~\ref{lem:weighted-gaussian-product}, on \(\mathcal G\), with conditional
probability at least \(1-4e^{-L}\),
\[
\|U\Gamma V^\top\|_{\mathrm{op}}
\le
C\left(\sqrt m\sqrt{n+L}+\sqrt{\frac mL}(n+L)\right).
\]
The same bound holds for \(V\Gamma U^\top\). Thus
\[
\|N_{\perp\perp}\|_{\mathrm{op}}
\le
C\left(\sqrt{\frac{n+L}{m}}+\frac{n+L}{\sqrt{mL}}\right)
\le
C\left(\sqrt{\frac dm}+\frac d{\sqrt{mL}}\right).
\]
Combining block estimates gives
\[
\|N(e_1)\|_{\mathrm{op}}
\le
C\left(\sqrt{\frac dm}+\frac d{\sqrt{mL}}\right).
\]

For the Frobenius bound, Lemma~\ref{lem:weighted-gaussian-product} gives
\[
\|U\Gamma V^\top\|_F^2\le Cn^2m,
\]
because \(\|\gamma\|_2^2\le Cm\), \(\|\gamma\|_\infty^2L\le Cm\), and
\(n=d-1\ge L\). Thus
\[
\|N_{\perp\perp}\|_F^2\le C\frac{d^2}{m}.
\]
The scalar and vector blocks satisfy
\[
|N_{11}|^2\le C\frac{d^2}{m},
\qquad
\|N_{1\perp}\|_2^2\le C\frac dm\le C\frac{d^2}{m}.
\]
Therefore
\[
\|N(e_1)\|_F^2\le C\frac{d^2}{m}.
\]

Finally,
\[
\operatorname{tr}N(e_1)
=
N_{11}+\operatorname{tr}N_{\perp\perp}
=
N_{11}+\frac1m\operatorname{tr}(U\Gamma V^\top).
\]
By Lemma~\ref{lem:weighted-gaussian-product},
\[
\frac1m|\operatorname{tr}(U\Gamma V^\top)|
\le
C\left(\sqrt{\frac{dL}{m}}+\sqrt{\frac Lm}\right)
\le
C\sqrt{\frac{dL}{m}}.
\]
The scalar term \(N_{11}\) is absorbed by the same bound. Hence
\[
|\operatorname{tr}N(e_1)|
\le
C\sqrt{\frac{dL}{m}}.
\]
This proves the fixed-direction estimate. Rotational invariance and the initial
union bound over the keys complete the proof.
\end{proof}

\begin{lemma}[Bilinear one-column quadratic-form parameters]
\label{lem:bilinear-column-parameters}
Assume the setting of Lemma~\ref{lem:bilinear-column-matrix-regularity}. Suppose
also that
\[
m\le \frac{d^2}{L},
\qquad
L^3\le c_0\sigma^2m^2,
\]
where
\[
\sigma^2=\Theta\left(\frac1{d^2}+\frac1m\right).
\]
For \(q\in\Sph\), define
\[
\tau(q)\defeq \frac{\operatorname{tr}A(q)}{d},
\qquad
B(q)\defeq A(q)-\tau(q)I_d,
\qquad
\beta(q)\defeq \tau(q)-\frac1d.
\]
Also define
\[
v(q)\defeq \frac{2\|B(q)\|_F^2}{d(d+2)}+\beta(q)^2,
\qquad
r(q)\defeq \frac{\|B(q)\|_{\mathrm{op}}}{d}+|\beta(q)|.
\]
If \(c_0>0\) is sufficiently small and \(C_0\) sufficiently large, then with
probability at least \(1-\delta/4\), simultaneously for every \(i\in[F]\),
\[
v(\vk_i)\le C\sigma^2,
\qquad
r(\vk_i)^2L^2\le C\sigma^2L.
\]
\end{lemma}

\begin{proof}
Since \(m\le d^2/L\),
\[
\frac1{d^2}\le \frac1{mL}\le \frac1m.
\]
Thus
\[
\sigma^2=\Theta\left(\frac1{d^2}+\frac1m\right)\asymp \frac1m.
\]
The assumption \(L^3\le c_0\sigma^2m^2\) implies \(m\ge C_*L^3\), after choosing
\(c_0>0\) sufficiently small. Hence
Lemma~\ref{lem:bilinear-column-matrix-regularity} applies. Also
\[
d^2\ge mL\ge C_*L^4,
\]
so \(d\ge cL^2\).

Work on the event from Lemma~\ref{lem:bilinear-column-matrix-regularity}. Fix
\(i\), write \(q=\vk_i\), \(A=A(q)\), and \(N=N(q)=A-qq^\top\). Since
\[
A=qq^\top+N,
\]
we have
\[
\operatorname{tr}A=1+\operatorname{tr}N,
\qquad
\beta(q)=\frac{\operatorname{tr}N}{d}.
\]
Using the trace bound,
\[
|\beta(q)|\le C\sqrt{\frac{L}{dm}}.
\]
Since \(d\ge L\),
\[
\beta(q)^2\le C\frac{L}{dm}\le C\frac1m\le C\sigma^2.
\]

Next,
\[
B(q)
=
\left(qq^\top-\frac1dI_d\right)
+
\left(N-\frac{\operatorname{tr}N}{d}I_d\right).
\]
The trace-removal map is an orthogonal projection in Frobenius norm, and
\[
\left\|qq^\top-\frac1dI_d\right\|_F^2
=
1-\frac1d
\le1.
\]
Therefore
\[
\|B(q)\|_F^2
\le
C\left(1+\|N\|_F^2\right)
\le
C\left(1+\frac{d^2}{m}\right).
\]
Hence
\[
\frac{2\|B(q)\|_F^2}{d(d+2)}
\le
C\left(\frac1{d^2}+\frac1m\right)
\le
C\sigma^2.
\]
Together with \(\beta(q)^2\le C\sigma^2\), this proves \(v(q)\le C\sigma^2\).

For \(r(q)\), use
\[
\|B(q)\|_{\mathrm{op}}
\le
1+\|N\|_{\mathrm{op}}+\frac{|\operatorname{tr}N|}{d}.
\]
Therefore
\[
\frac{\|B(q)\|_{\mathrm{op}}}{d}
\le
C\left(
\frac1d+
\frac1{\sqrt{dm}}+
\frac1{\sqrt{mL}}+
\sqrt{\frac{L}{d^3m}}
\right).
\]
Since \(d\ge L^2\), all terms except \(1/d\) are dominated by
\(C/\sqrt{mL}\). Also
\[
|\beta(q)|
\le
C\sqrt{\frac{L}{dm}}
\le
C\frac1{\sqrt{mL}}.
\]
Thus
\[
r(q)
\le
C\left(\frac1d+\frac1{\sqrt{mL}}\right).
\]
Squaring and multiplying by \(L^2\),
\[
r(q)^2L^2
\le
C\left(\frac{L^2}{d^2}+\frac{L}{m}\right).
\]
Since \(m\le d^2/L\),
\[
\frac{L^2}{d^2}\le \frac{L}{m}.
\]
Therefore
\[
r(q)^2L^2\le C\frac{L}{m}\le C\sigma^2L.
\]
The estimates hold simultaneously for all \(q=\vk_i\).
\end{proof}
